% paper_18_exchange_constants.tex — GeoVac Paper 18
% Spectral-Geometric Exchange Constants: Projecting Discrete
% Spectra onto Continuous Manifolds
% Status: Draft, March 29, 2026
% Category: Observations (papers/observations/)
\documentclass[aps,pra,twocolumn,superscriptaddress,longbibliography]{revtex4-2}

\usepackage{amsmath,amssymb,amsthm,graphicx,xcolor,bm,braket,dcolumn,booktabs}
\usepackage{hyperref}

\newtheorem{theorem}{Theorem}
\newtheorem{observation}{Observation}
\newtheorem{conjecture}{Conjecture}

\begin{document}

\title{Spectral-Geometric Exchange Constants:\\
Projecting Discrete Spectra onto Continuous Manifolds}
\thanks{Paper 18 in the Geometric Vacuum series}

\author{J.~Loutey}
\affiliation{Independent Researcher, Kent, Washington}

\date{March 29, 2026}

%% ========================================================================
%% ABSTRACT
%% ========================================================================

\begin{abstract}
Compact Lie groups have discrete unitary representations
(Peter--Weyl theorem); the GeoVac graph exploits this discreteness
directly.  We show that the transcendental conversion factors
appearing throughout the GeoVac hierarchy---$\kappa = -1/16$
($S^3 \to \mathbb{R}^3$), the Stieltjes seed $e^a E_1(a)$
(Laguerre basis $\to$ prolate spheroid), the adiabatic eigenvalue
$\mu(R)$ ($S^3 \times S^3 \to$ hyperspherical)---are manifestations
of a single structural pattern: they arise when discrete
representations of compact groups are projected onto non-compact
coordinate systems used for measurement.  The taxonomy classifies
these factors by the type of compactness mismatch they mediate,
and connects them to the Weyl--Selberg exchange constants of
spectral geometry.  This identification explains why the constants
are transcendental (compact$\to$non-compact projection introduces
$\pi^{d/2}$), why their complexity increases with the geometry,
and why they vanish when the computation remains on the graph.
The falsifiable claim: the transcendental content of any GeoVac
observable is determined entirely by the type of projection
linking the compact discrete description to the non-compact
coordinate system.  We note that the combination rule
$K = \pi(B + F - \Delta)$ in the conjectural fine structure
constant formula (Paper~2) has the same structural character:
$\pi$ multiplying a sum of discrete representation-theoretic
data and continuous spectral invariants.
\end{abstract}

\maketitle

%% ========================================================================
%% I. INTRODUCTION
%% ========================================================================

\section{Introduction}
\label{sec:intro}

The Geometric Vacuum program constructs discrete lattice
Hamiltonians by identifying the natural geometry of each quantum
system and discretizing it as a sparse
graph~\cite{loutey_paper0,loutey_paper1,loutey_paper7}.  A
persistent observation across the hierarchy of natural geometries
is that the discrete algebraic structures---integer eigenvalues,
finite recurrence relations, Gaunt integrals built from Wigner
$3j$ symbols---are entirely self-contained, but that their
connection to physical observables in continuous coordinate
systems always introduces transcendental constants.

The thesis of this paper is that these transcendental conversion
factors are consequences of \emph{compactness}.  Compact Lie groups
have discrete unitary representations (the Peter--Weyl
theorem~\cite{peter_weyl1927}); the GeoVac graph exploits this
discreteness directly.  Transcendental numbers enter when---and only
when---those discrete representations are translated into non-compact
coordinate systems used for measurement.  The taxonomy developed below
classifies the resulting conversion factors by the type of
compact-to-non-compact projection that produces them.  This yields a
falsifiable structural claim: the transcendental content of any GeoVac
observable is determined entirely by the type of projection linking
the compact discrete description to the non-compact coordinate system
(Claim~4, Sec.~\ref{sec:observable_classification}).

The simplest instance is already visible in Paper~0's packing
construction~\cite{loutey_paper0}.  Axiom~2 (geometric indifference)
forces states onto isotropic shells---circles, i.e.\ copies of~$S^1$.
The shell is compact; its circumference $2\pi r$ introduces $\pi$
into the area formula $\sigma_0 = \pi d_0^2/2$ and hence into the
state count $N_k = A_k / \sigma_0 = 2(2k-1)$.  The integer state count
is a consequence of compact topology; $\pi$ is the price of expressing
that count as a ratio of continuous areas.  This pattern---compact
geometry producing integer quantum numbers, with $\pi$ entering
through projection onto non-compact coordinates---recurs at every level
of the hierarchy.

The mathematical framework is not new: the conversion factors are
instances of Weyl's law~\cite{weyl1911} and the Selberg trace
formula~\cite{selberg1956}.  What is new is their systematic
identification across the GeoVac hierarchy, and the observation
that the conjectural fine structure constant formula (Paper~2)
has the same structural character.


%% ========================================================================
%% II. CATALOG OF EXCHANGE CONSTANTS
%% ========================================================================

\section{Catalog of Exchange Constants}
\label{sec:catalog}

\subsection{Level 1: $\kappa = -1/16$ ($S^3 \to \mathbb{R}^3$)}
\label{sec:kappa}

The compactness of $S^3$ is what makes the graph eigenvalues discrete
integers; $\kappa$ converts this discrete spectrum to the non-compact
$\mathbb{R}^3$ energy scale.

The graph Laplacian on $S^3$ produces pure integer eigenvalues
\begin{equation}
  \lambda_n = -(n^2 - 1), \qquad n = 1, 2, 3, \ldots
  \label{eq:s3_eigenvalues}
\end{equation}
with degeneracy $g_n = n^2$~\cite{loutey_paper7}.  The physical
hydrogen energy levels $E_n = -1/(2n^2)$ are related by the
energy-shell constraint $p_0^2 = -2E$, which acts as a
stereographic focal length projecting $S^3$ onto flat
$\mathbb{R}^3$.  The universal kinetic scale factor
\begin{equation}
  \kappa = -\frac{1}{16}
  \label{eq:kappa}
\end{equation}
calibrates the ground state energy:
$\kappa \cdot \lambda_{\max} = E_H = -1/2$~Ha, where
$\lambda_{\max} \to 8$ is the maximum graph Laplacian eigenvalue
in the continuum limit~\cite{loutey_paper0}.  Per-shell energies
$E_n = -Z^2/(2n^2)$ emerge from the full graph Hamiltonian
$H = \kappa \mathcal{L} + W$ (with node weights $W_{nn} = -Z/n^2$),
not from $\kappa$ alone.

The origin of $\kappa$ is well understood geometrically: it is
the Jacobian of the stereographic projection composed with the
conformal factor of the Fock map.  Specifically, if $\Omega$
is the conformal factor $\Omega = (1 + p^2/p_0^2)^{-1}$, then
the Laplace--Beltrami operator on $S^3$ and the flat-space
kinetic operator are related by
$\nabla^2_{S^3} = \Omega^{-2}[\nabla^2_{\mathrm{flat}} + \text{lower order}]$,
and $\kappa$ absorbs the accumulated conformal rescaling and
the $p_0$ energy-shell normalization~\cite{loutey_paper7}.

The key property of $\kappa$ is that it is \textit{unnecessary
on the graph}.  The graph Laplacian eigenvalue problem $L\psi =
\lambda\psi$ is solved entirely in terms of integers.  The
constant $\kappa$ appears only when the result is translated
into atomic units (Hartrees, bohrs)---i.e., when the discrete
spectrum is projected onto a continuous coordinate system.

\subsection{Level 2: $e^a E_1(a)$ (Laguerre basis $\to$ prolate spheroid)}
\label{sec:stieltjes}

The Laguerre spectral basis inherits the compact angular structure of
$S^3$; the Stieltjes seed $e^a E_1(a)$ enters when this basis is
embedded in the non-compact prolate spheroidal coordinate, whose
curvature introduces a pole absent from the compact description.

The prolate spheroidal radial solver uses a Laguerre spectral
basis $\varphi_n(\xi) = e^{-\alpha(\xi-1)} L_n(2\alpha(\xi-1))$
with matrix elements computed from three-term
recurrences~\cite{loutey_paper11}.  For $\sigma$ states ($m = 0$),
all matrix elements are purely algebraic: the Laguerre moment
matrices $M_k[i,j] = \int_0^\infty x^k L_i(x) L_j(x) e^{-x}
dx$ are band-structured matrices with rational entries, built
from the recurrence $x L_n = -(n+1)L_{n+1} + (2n+1)L_n -
nL_{n-1}$.

For $m \neq 0$ states, the centrifugal term $m^2/(\xi^2 - 1)$
introduces a $1/x$ singularity in the Laguerre variable
$x = 2\alpha(\xi - 1)$.  The associated Laguerre basis
$L_n^{|m|}(x)$ with weight $x^{|m|} e^{-x}$ absorbs this
singularity by construction.  The centrifugal matrix elements
decompose via partial-fraction expansion:
\begin{equation}
  \frac{m^2}{\xi^2 - 1}
  = 4\alpha^2 m^2 \left[\frac{1}{x}
    - \frac{1}{x + 4\beta}\right],
  \label{eq:partial_fraction}
\end{equation}
where $\beta = \alpha(\xi_0 - 1)$ is the Laguerre scale offset.
The two terms have fundamentally different algebraic character:
\begin{itemize}
  \item The $1/x$ term produces a lowered moment
    $M_{-1}[i,j] = \int_0^\infty x^{|m|-1} L_i^{|m|}(x)
    L_j^{|m|}(x)\,e^{-x}\,dx$, which is fully algebraic via
    the DLMF 18.9.13 summation
    identity~\cite{dlmf_laguerre}.
  \item The $1/(x + 4\beta)$ term produces a Stieltjes matrix
    $J[i,j] = \int_0^\infty \frac{x^{|m|}}{x + a}\,
    L_i^{|m|}(x) L_j^{|m|}(x)\,e^{-x}\,dx$, where $a = 4\beta$.
    This matrix satisfies a three-term recurrence in the
    row/column indices, seeded by the single transcendental
\end{itemize}
\begin{equation}
  S_0^{(\alpha)}(a) = \Gamma(\alpha+1)\,e^a\,E_{\alpha+1}(a),
  \label{eq:stieltjes_seed}
\end{equation}
where $E_n(a) = \int_1^\infty t^{-n} e^{-at}\,dt$ is the
generalized exponential integral and $a = 4\beta$ is the
Laguerre scale parameter~\cite{loutey_track_j}.  For integer
$\alpha = |m|$, all higher $E_n$ reduce to $E_1(a)$ via a
finite rational recurrence, so the entire matrix is built from
algebraic operations plus one transcendental constant.

The computational cost of the transcendental seed is negligible:
one call to $E_1(a)$, which achieves machine precision for
$a \ge 4$ (covering all physical parameters).  The seed then
propagates through an $O(N^2)$ algebraic recurrence that produces
all matrix elements.  The continuous mathematics has been
compressed from a \textit{function} (the quadrature integrand
over $[0,\infty)$) to a \textit{constant} (the Stieltjes seed),
with the entire function-level information reconstructed
algebraically from that constant.

An unexpected computational bonus: the associated Laguerre basis
converges \textit{faster} than the ordinary Laguerre basis for
$m \neq 0$ states.  For $|m| = 1$ ($\pi$ states), the associated
basis is stable by $N = 10$, compared to non-monotonic convergence
at $N \approx 20$ for the ordinary basis.  For $|m| = 2$
($\delta$ states), algebraic-quadrature agreement reaches
$\sim 5 \times 10^{-9}$ (machine precision), with PES shape
(equilibrium geometry) preserved exactly~\cite{loutey_track_j}.
This suggests that identifying the exchange constant and absorbing
it into the basis does not merely explain the transcendence---it
\textit{improves the spectral method}, because the adapted basis
eliminates the curvature mismatch that the exchange constant
mediates.

The Stieltjes seed $e^a E_1(a)$ is the Laplace transform of the
rational function $1/(1+t)$ evaluated at $a$.  It encodes how
the Laguerre inner product (defined on the flat half-line
$[0, \infty)$ with exponential weight) ``sees'' a shifted pole
at $x = -a$ arising from the curvature of the prolate spheroidal
coordinate system.  On the graph, there is no pole; the
centrifugal barrier is a selection rule encoded in the quantum
numbers.  The transcendental appears only when this selection
rule is represented as a rational function in a continuous
coordinate.

\subsection{Level 3: $\mu(R)$ ($S^3 \times S^3 \to$ hyperspherical)}
\label{sec:mu}

On the compact product space $S^3 \times S^3$, the two-electron
Hamiltonian has purely algebraic matrix elements; the adiabatic
eigenvalue $\mu(R)$ enters only when this compact description is
reparameterized into the non-compact hyperspherical coordinate~$R$.

The two-electron problem on $S^3 \times S^3$ is fully algebraic:
the Hamiltonian matrix has entries given by Gaunt integrals
(finite sums of Wigner $3j$ symbols with rational coefficients),
and the eigenvalues are roots of a polynomial with algebraic
coefficients.  The Full CI solver exploits this, achieving
0.35\% error for helium with no transcendentals
anywhere~\cite{loutey_fci}.

The hyperspherical reparameterization introduces coordinates
$(R, \alpha)$ where $R = \sqrt{r_1^2 + r_2^2}$ is the
hyperradius and $\alpha = \arctan(r_2/r_1)$ is the hyperangle.
The angular eigenvalues $\mu(R)$ depend parametrically on $R$
because the effective angular Hamiltonian is
$H(R) = H_0 + R \cdot V^{\mathrm{coupling}}$~\cite{loutey_paper13}.

Each perturbation coefficient $a_k$ in the expansion
$\mu(R) = \mu^{(0)} + a_1 R + a_2 R^2 + \cdots$ has a closed
form: the $a_k$ are finite sums of algebraic numbers (square
roots of rationals from Clebsch--Gordan coefficients) times
powers of $\pi^{-1}$ (from normalization over the angular domain
$[0, \pi/2]$), arising from partial harmonic sums over
$\mathrm{SO}(6)$ Casimir eigenstates~\cite{loutey_paper13}.
For example, $a_1(n{=}1, Z) = (8/3\pi)(\sqrt{2} - 4Z)$.
However, the full function $\mu(R)$ is not transcendental in the
sense originally claimed.  The angular Hamiltonian
$H(R) = H_0 + R \cdot V_C$ is a linear matrix pencil, so its
eigenvalues satisfy the characteristic polynomial
\begin{equation}
  P(R, \mu) \;=\; \det\!\bigl(H_0 + R\,V_C - \mu\,I\bigr) = 0,
  \label{eq:char_poly}
\end{equation}
which is a polynomial of degree $d$ in both $R$ and $\mu$ (where
$d$ is the matrix dimension at a given truncation).  The coefficients
of $P$ lie in $\mathbb{Q}(\pi, \sqrt{2})$, since $H_0$ has integer
entries (SO(6) Casimir eigenvalues) and $V_C$ has entries involving
$\pi^{-1}$ and $\sqrt{2}/\pi$ from the angular integrals.  This
makes $\mu(R)$ an \emph{algebraic} function of $R$ over the
coefficient ring $\mathbb{Q}(\pi, \sqrt{2})$---it is a root of
a polynomial, not a transcendental
function~\cite{loutey_track_p1}.

The transcendental content of $\mu(R)$ is therefore
\emph{inherited} from the coupling matrix $V_C$, which contains
upstream exchange constants of embedding type ($\pi^{-1}$ from
angular normalization, $\sqrt{2}$ from Gaunt-type integrals).
The $R$-dependence itself introduces no new transcendentals:
it is purely algebraic.

The failure of the Taylor/Pad\'e expansion at large $R$
(Track~G~\cite{loutey_track_g}) is reinterpreted as a
\emph{branch point} phenomenon, not evidence of transcendence.
The algebraic curve $P(R, \mu) = 0$ has branch points where
eigenvalue sheets meet---these are the zeros of the discriminant
$\Delta(R) = \mathrm{disc}_\mu(P)$.  The Taylor series around
$R = 0$ has a convergence radius set by the nearest branch point
in the complex $R$-plane.  Pad\'e approximants, being rational
functions, also cannot cross branch points where $\mu(R)$ has
square-root type singularities.  But the implicit relation
$P(R, \mu) = 0$ is exact for \emph{all} $R$, including $R > 5$
where the series diverges.  This has been verified computationally
to machine precision at $R = 50$~bohr~\cite{loutey_track_p1}.

The $O(R) \to O(R^2)$ crossover, which was previously interpreted
as a regime change requiring transcendental interpolation, is
simply the behavior of an algebraic function whose dominant term
changes character: at small $R$, the linear term $V_C \cdot R$
dominates; at large $R$, higher-order terms in the characteristic
polynomial dominate.  The crossover is smooth along the algebraic
curve, even though it is singular from the perspective of a Taylor
expansion centered at $R = 0$.

The crucial point remains that $\mu(R)$ is not intrinsic to the
two-electron problem.  It is an artifact of the hyperspherical
coordinate system.  The $S^3 \times S^3$ product space does not
require $\mu(R)$; it only appears when one reparameterizes for
computational efficiency.  The exchange rate between the discrete
product-space description and the continuous hyperspherical
description is an algebraic function of $R$, with transcendental
content inherited from the angular coupling integrals.


\subsection{Level 4: $\mu(\rho,R)$ (product space $\to$ molecule-frame hyperspherical)}
\label{sec:mu_level4}

At Level~4 the compact product basis (prolate spheroidal CI) is fully
algebraic but structurally incomplete; the non-compact mol-frame
hyperspherical reparameterization introduces $\mu(\rho,R)$ as the
irreducible cost of accessing three-body coalescence physics that the
compact description cannot represent.

Level~4 combines the two-center complication of Level~2 with the
two-electron complication of Level~3.  Two approaches exist, with
fundamentally different algebraic character.

The product-space approach (prolate spheroidal CI, Paper~12) is
fully algebraic: all electron-electron matrix elements are computed
from Neumann expansion recurrence relations with zero spatial
quadrature.  The Laguerre radial and Legendre angular matrix
elements are likewise algebraic.  No transcendental constant
appears anywhere in the computation.

However, a systematic convergence study (Track~M) demonstrates
that this algebraic approach saturates at $\sim$92.5\% of the
exact dissociation energy $D_e$~\cite{loutey_track_m}.  The
angular basis convergence is rapid---92.3\% at $l_{\max} = 2$
(46 basis functions), 92.5\% at $l_{\max} = 4$ (114 basis
functions)---but successive improvements decay by a factor of
200 from $l = 1 \to 2$ to $l = 2 \to 3$, with the $l = 3 \to 4$
increment contributing only 0.04\%.  Extrapolation in the slow
regime predicts that 5\% error would require $l_{\max} \sim 45$
and 1\% error $l_{\max} \sim 113$, both impractical.  The
effective ceiling is 92.5\% $D_e$.

The smoking gun is explicit correlation: 9~basis functions with
$r_{12}^2$ factors (Hylleraas $p = 2$) achieve 94.7\% $D_e$,
exceeding the 114-function prolate CI plateau.  The residual
7.5\% matches Paper~12's cusp diagnosis (7.6\%): the
electron-electron cusp ($1/r_{12}$ singularity) is a completeness
deficiency of the product basis, not a convergence issue.
The transcorrelated approach (Paper~14, Sec.~IV; Paper~15,
Sec.~VI.J) removes this singularity via a Jastrow similarity
transformation, converting $1/r_{12}$ from a hard embedding
constant to a smooth gradient operator at the cost of
non-Hermiticity (Sec.~\ref{sec:taxonomy}).

The mol-frame hyperspherical reparameterization (Paper~15)
resolves this by introducing coordinates $(\rho, R, \alpha)$
where $\rho$ is the electron hyperradius and $\alpha$ the
hyperangle.  The angular eigenvalues $\mu(\rho, R)$ are now a
transcendental function of \textit{two} continuous parameters,
reflecting the simultaneous three-body coalescence in coordinate
and configuration space.  This solver achieves 94.1\% $D_e$,
surpassing the prolate CI ceiling.

The distinction between Levels~3 and 4 is qualitative.  At
Level~3, the $S^3 \times S^3$ FCI converges (albeit slowly)---the
exchange constant $\mu(R)$ is a computational convenience that
accelerates convergence but is not fundamentally necessary.  At
Level~4, the prolate product-space CI \textit{saturates}: the
exchange constant $\mu(\rho, R)$ is irreducible because the
product basis is structurally blind to three-body coalescence.
The reparameterization does not merely accelerate convergence;
it accesses physics that the algebraic product space cannot
represent at any practical basis size.


%% ========================================================================
%% III. THE PATTERN
%% ========================================================================

\section{The Pattern: Weyl--Selberg Exchange}
\label{sec:pattern}

\subsection{Weyl's law}

The mathematical content of the observations in
Sec.~\ref{sec:catalog} is an instance of Weyl's law, which
states that the eigenvalue counting function $N(\lambda)$ of
the Laplace--Beltrami operator on a compact Riemannian manifold
of dimension $d$ satisfies~\cite{weyl1911}
\begin{equation}
  N(\lambda) \sim \frac{\omega_d}{(2\pi)^d}\,
  \mathrm{vol}(\Omega)\,\lambda^{d/2}
  \qquad (\lambda \to \infty),
  \label{eq:weyl}
\end{equation}
where $\omega_d = \pi^{d/2}/\Gamma(d/2+1)$ is the volume of
the unit ball in $\mathbb{R}^d$.

Weyl's law relates two fundamentally different descriptions
of the same manifold:
\begin{itemize}
  \item The \textbf{spectral side}: $N(\lambda)$, a discrete
    staircase function counting integer-labeled eigenvalues.
  \item The \textbf{geometric side}: $\mathrm{vol}(\Omega)$, a
    continuous quantity measured by integration over the manifold.
\end{itemize}
The conversion factor $\omega_d / (2\pi)^d$ involves
$\pi^{d/2}$---a transcendental that depends only on the
dimension.  This factor is the prototypical
\textit{spectral-geometric exchange constant}: it converts
between counting (discrete) and measuring (continuous).

For the specific manifolds in the GeoVac hierarchy, the ball
volumes $\omega_d$ carry $\pi^{\lfloor d/2 \rfloor}$:
\begin{itemize}
  \item $S^1$ ($d = 1$): $\omega_1 = 2$,
    Weyl constant $\omega_1/(2\pi)^1 = 1/\pi$.
  \item $S^2$ ($d = 2$): $\omega_2 = \pi$,
    Weyl constant $\omega_2/(2\pi)^2 = 1/(4\pi)$.
  \item $S^3$ ($d = 3$): $\omega_3 = 4\pi/3$,
    Weyl constant $\omega_3/(2\pi)^3 = 1/(6\pi^2)$.
  \item $S^5$ ($d = 5$): $\omega_5 = 8\pi^2/15$,
    Weyl constant $\omega_5/(2\pi)^5 = 1/(60\pi^3)$.
\end{itemize}
In each case, the Weyl constant involves inverse powers of $\pi$
(specifically $\pi^{-\lceil d/2 \rceil}$), reflecting the
transcendental cost of converting a discrete eigenvalue count
into a continuous volume measure.

\subsection{Compactness as the source of discreteness}

The Peter--Weyl theorem states that the unitary representations of a
compact Lie group form a discrete, countable, orthogonal decomposition
of $L^2(G)$~\cite{peter_weyl1927}.  This is the structural reason that
quantum numbers are integers: $\mathrm{SO}(4)$ is compact, so its
irreducible representations are labeled by discrete indices $(n, l, m)$
with finite multiplicities $g_n = n^2$.  The graph Laplacian on $S^3$
inherits this discreteness directly---its eigenvalues $\lambda_n =
-(n^2-1)$ are integers because the underlying group is compact.

Exchange constants arise when these discrete representations are
expressed in non-compact coordinate systems.  The stereographic
projection $S^3 \to \mathbb{R}^3$ maps the compact manifold (finite
volume $2\pi^2$) onto flat space (infinite volume), introducing the
conformal factor $\Omega = (1 + p^2/p_0^2)^{-1}$ and the calibration
constant $\kappa$.  More generally, every exchange constant in
the taxonomy (Sec.~\ref{sec:taxonomy}) mediates a
compact$\to$non-compact translation: $e^a E_1(a)$ converts compact
angular structure into a non-compact radial coordinate with a pole;
$\mu(R)$ converts the compact product space $S^3 \times S^3$ into
non-compact hyperspherical coordinates.  The transcendental content
is the toll for crossing the compactness boundary.

\subsection{$\pi$ as the founding example}

The simplest spectral-geometric exchange constant is $\pi$ itself.
The Laplacian on $S^1$ (the circle) has eigenvalues $n^2$, counted
by integers; the continuous circumference is $2\pi$.  The
conversion between the discrete eigenvalue count and the continuous
measure is mediated by $\pi$.  More generally, the volume of the
unit $d$-sphere involves $\pi^{d/2}$, and the Weyl constant
$\omega_d / (2\pi)^d$ carries exactly this factor.  Every exchange
constant in the GeoVac hierarchy is a descendant of this primordial
example: $\kappa$ inherits from the $S^3$ Weyl constant, the
Stieltjes seed involves the exponential integral (a function whose
Taylor coefficients involve the Euler--Mascheroni constant, itself
a limit of harmonic sums weighted by $\pi$-dependent terms), and
$\mu(R)$ is an eigenvalue of a Hamiltonian on $S^5$ whose Weyl
asymptotics carry $\pi^{5/2}$.

\begin{observation}[The graph is $\pi$-free]
\label{obs:pi-free}
Every algebraic operation in the GeoVac graph construction
produces integers or rationals:
\begin{itemize}
  \item Eigenvalues of the graph Laplacian: $\lambda_n = -(n^2 - 1)$,
    integer for all $n$.
  \item Degeneracies: $g_n = n^2$, integer.
  \item Gaunt integrals (Legendre normalization,
    $\int P_{l_1} P_k P_{l_2}\,dx$): ratios of factorials via
    Wigner $3j$ symbols, rational.
  \item Selection rules: combinatorial (triangle inequality on angular
    momentum labels).
  \item Clebsch--Gordan coefficients: algebraic numbers (square roots
    of rationals).
  \item Casimir eigenvalues: $l(l+1)$, $\nu(\nu+4)/2$, integer or
    half-integer for all levels.
\end{itemize}
Transcendental numbers ($\pi$, $e^a E_1(a)$, $\zeta(2)$) enter the
framework exclusively through normalization to continuous measures:
spherical harmonic normalization ($\int |Y|^2\,d\Omega = 1$ with
$d\Omega$ carrying $4\pi$), manifold volumes
($\mathrm{vol}(S^3) = 2\pi^2$), and spectral zeta functions
($\zeta_{S^1}(1) = \pi^2/6$).  The exchange constants cataloged in
this paper are precisely the $\pi$-injection points---the places
where the discrete graph contacts continuous geometry.
\end{observation}

\subsection{The Selberg trace formula}

Weyl's law is asymptotic.  The Selberg trace
formula~\cite{selberg1956} provides the exact identity:
\begin{equation}
  \sum_j h(r_j)
  = \frac{\mathrm{Area}(\Gamma \backslash \mathcal{H})}{4\pi}
    \int_{-\infty}^{\infty} r\,\tanh(\pi r)\,h(r)\,dr
    + \sum_{\{\gamma\}} \cdots
  \label{eq:selberg}
\end{equation}
where the left side sums over the discrete spectrum, the first
term on the right involves the continuous volume (with a factor
of $\pi$), and subsequent terms sum over conjugacy classes
(geodesics).  The Selberg trace formula is the bridge between
the spectral world and the geometric world, and the
transcendental factors ($\pi$, $\tanh$, logarithms of
algebraic numbers) are the toll for crossing.

When the quotient $\Gamma \backslash G$ is compact, the spectrum
is purely discrete and the trace formula involves only algebraic
data plus $\pi$-dependent conversion factors.  When the quotient
is non-compact, a continuous spectrum appears and the exchange
involves Eisenstein series---transcendental
functions~\cite{selberg1956}.  This exactly parallels the
GeoVac hierarchy: Level~1 (compact $S^3$) has a purely algebraic
spectrum plus a constant ($\kappa$); Level~3 ($S^5$ with $\mathrm{SO}(6) \to \mathrm{SO}(4)$
symmetry crossover) has an algebraic function $\mu(R)$ with
transcendental coefficients (Sec.~\ref{sec:algebraic_curve}).


%% ========================================================================
%% IV. TAXONOMY
%% ========================================================================

\section{Taxonomy of Exchange Constants}
\label{sec:taxonomy}

The following taxonomy classifies exchange constants by the type of
compactness mismatch they mediate.  The hierarchy---intrinsic,
calibration, embedding, flow, composition---can be read as increasing
departure from compactness: intrinsic constants convert between
compact manifolds of different dimension; calibration constants convert
between a compact manifold and flat space; embedding constants arise
from coordinate singularities in the non-compact description;
flow functions encode parameter-dependent symmetry breaking; and
composition constants measure the cost of factorizing a single compact
space into multiple incompatible coordinate systems.

Based on the catalog and the Weyl--Selberg identification, we
propose the following taxonomy:

\begin{table}[h]
\caption{Spectral-geometric exchange constants in the GeoVac
  hierarchy.  ``Source'' indicates the discrete structure;
  ``Target'' indicates the continuous coordinate system.
  ``Type'' classifies the exchange constant by what determines it.
  ``Algebraic (implicit)'' indicates that the function satisfies a
  polynomial equation $P(R,\mu) = 0$ with coefficients inherited from
  upstream exchange constants (Sec.~\ref{sec:algebraic_curve}).
  ``Piecewise'' indicates that the angular Hamiltonian has
  piecewise-smooth parameter dependence (from the split-region
  Legendre expansion), preventing a global characteristic polynomial.
  The function is piecewise-algebraic within each region.}
\label{tab:taxonomy}
\begin{ruledtabular}
\begin{tabular}{ccccc}
  Level & Source & Target & Constant & Type \\
  \hline
  --- & $S^1$ lattice & $\mathbb{R}^1$ & $\pi$ & intrinsic \\
  1 & $S^3$ graph & $\mathbb{R}^3$ & $\kappa = -1/16$ & calibration \\
  2 ($m{=}0$) & Laguerre rec. & prolate sph. & (none) & algebraic \\
  2 ($m{\neq}0$) & assoc. Laguerre & prolate sph. & $e^a E_1(a)$
    & embedding \\
  3 (FCI) & $S^3 {\times} S^3$ & (none) & (none) & algebraic \\
  3 (hyper.) & $\mathrm{SO}(6)$ reps & $S^5(R)$ param.
    & $\mu(R)$ & alg.\ (implicit) \\
  4 (prolate CI) & prolate product & (none) & (none) & algebraic \\
  4 (hyper.) & $\mathrm{SO}(6)$ reps & $(\rho,R)$ param.
    & $\mu(\rho,R)$ & piecewise \\
  5 (composed) & Level 3+4 fibers & full $N$-electron
    & PK pseudopot. & composition \\
  5 (balanced) & hydrogenic $\times$ multipole & cross-center $V_{ne}$
    & $\gamma(k,\alpha R)$ & algebraic \\
\end{tabular}
\end{ruledtabular}
\end{table}

The taxonomy reveals a hierarchy of exchange types:
\begin{enumerate}
  \item \textbf{Algebraic} (no exchange constant): When the
    computation stays on the graph or uses a basis whose inner
    product matches the graph structure, no transcendental
    appears.  Examples: Level~1 on $S^3$, Level~2 $\sigma$
    states, Level~3 FCI.

  \item \textbf{Intrinsic constants} ($\pi^{d/2}$): Determined
    by the manifold's dimension and volume alone, independent of
    any discretization or basis choice.  These are the Weyl
    constants---universal volume-to-eigenvalue conversion factors.
    Example: $\pi$ itself, for $S^1 \to \mathbb{R}$.

  \item \textbf{Calibration constants}: Determined by the
    discrete graph construction, not by intrinsic manifold
    geometry.  Depends on lattice connectivity.  Example:
    $\kappa = -E_H/\lambda_{\max} = -1/16$, where
    $\lambda_{\max} = 2 d_{\max} = 8$ depends on the maximum
    vertex degree $d_{\max} = 4$ of the GeoVac graph
    (Sec.~\ref{sec:kappa_derivation}).  A different graph with
    $d_{\max} = 6$ would give $\kappa = -1/24$.

  \item \textbf{Embedding constants}: Single transcendental
    numbers seeding algebraic recurrences.  Determined by the
    interaction between a spectral basis and a coordinate
    singularity.  Examples: $e^a E_1(a)$ for the prolate
    centrifugal pole; the $1/r_{12}$ electron-electron cusp,
    which is the embedding exchange constant for the two-electron
    coalescence in any product basis.
    The cusp embedding constant is unique in being \emph{removable}
    by a similarity transformation: the transcorrelated (TC)
    Jastrow factor $J = \tfrac{1}{2} r_{12}$ satisfies
    $e^{-J} (1/r_{12})\, e^{J} = \text{smooth}$, converting
    the algebraic partial-wave convergence $1/(l+\tfrac{1}{2})^4$
    to exponential.  This makes $1/r_{12}$ a ``soft'' embedding
    constant---present in any product-space representation but
    eliminable by a non-unitary change of basis---in contrast to
    $e^a E_1(a)$, which is structurally irreducible.

    The $1/r_{12}$ classification is \emph{universal across second-order
    central-field operators}: Sprint~2 Track~DC-A (April~2026) extended
    the analysis to the two-electron Dirac-Coulomb Hamiltonian and
    proved, via the Pauli-reduced LL block (Kutzelnigg 1984, 1988), that
    the cusp condition $(\partial_{r_{12}} \Psi_{LL})|_{r_{12}=0^+} = \tfrac12 [1+O(\alpha^2)] \Psi_{LL}(0)$
    survives relativistic corrections with only an
    $O((Z\alpha)^2)$ amplitude rescaling---the Kutzelnigg-Morgan 1992
    Theorem~3 Sobolev-regularity class is preserved.  The partial-wave
    energy convergence exponent is $l^{-4}$ for singlet (Schwartz 1962)
    and $l^{-6}$ for triplet (Pack-Brown 1966) in both Schrödinger-Coulomb
    and Dirac-Coulomb.  Numerical verification at Z=4 (Be$^{2+}$) via
    Track~DC-B: fitted exponent matches to $\Delta p = +0.019$ within
    fit-noise.  The Dirac program therefore does \emph{not} regularize
    the e-e cusp; full QED (vacuum polarization modifying the photon
    propagator at distances $\lesssim 1/m_e c$) would be required to
    soften the $1/r_{12}$ kernel itself.

    A \emph{distributional} sub-category of embedding constants appears
    in Sprint~2 Track~BR-B.  The Breit spin-spin interaction contains a
    $1/r_{12}^3$ kernel whose bare partial-wave expansion
    \[
      \frac{1}{r_{12}^3} = \sum_l (2l+1)\, K_l(r_<, r_>)\, P_l(\cos\theta_{12}),
      \qquad K_l(r_<, r_>) = \frac{r_<^l}{r_>^{l+1}\,(r_>^2 - r_<^2)}
    \]
    has a \emph{logarithmic divergence} at coalescence $r_< \to r_>$
    for every $l$---the operator is not in $L^1_{\rm loc}$ and its
    matrix elements are undefined as ordinary integrals.  The Breit-Pauli
    tensor projection (transverse $\sigma_1\!\cdot\!\sigma_2 - 3(\sigma_1\!\cdot\!\hat r_{12})(\sigma_2\!\cdot\!\hat r_{12})$)
    regularizes this by eliminating the longitudinal component, leaving
    a smooth kernel $K_l^{BP}(r_<, r_>) = r_<^l / r_>^{l+3}$ whose matrix
    elements are exact rationals for same-$n$ orbitals (e.g.\
    $R_{BP}^0(1s,1s;1s,1s) = 0$, $R_{BP}^2(2p,2p;2p,2p) = 7/768$,
    $R_{BP}^1(2s,2p;2s,2p) = 1/256$ at $Z=1$) with $Z^3$ scaling.
    For \emph{different-$n$} orbitals, a new transcendental appears:
    $M_{\rm ret}^2(1s,2p;1s,2p) = 21344/729 - (10240/243)\log(2)$,
    the first documented example of a $\log$-embedding constant in the
    GeoVac framework.  The Breit-Pauli regularization is thus both
    \emph{distributional} (the bare operator requires regularization
    before having matrix elements) and \emph{multi-scale} (cross-$n$
    integrals introduce $\log$ content absent from same-$n$).

    A graph-theoretic characterization of $1/r_{12}$ on the $S^3$
    pair-state lattice (v2.9.1) makes the embedding nature explicit.
    Building the graph whose nodes are singlet pair-states
    $|(n_1 l_1 m_1)(n_2 l_2 m_2)\rangle$ and whose edges are the exact
    Slater matrix elements $\langle ab | 1/r_{12} | cd \rangle$, one
    finds: (i)~diagonal elements are exact rationals
    ($\langle 1s\,1s | V | 1s\,1s \rangle = 5/8$,
    $\langle 2s\,2s | V | 2s\,2s \rangle = 77/512$, etc.) but cross-shell
    denominators mix $2^k$ and $3^j$ factors, so the spectrum of $V$
    is not rational---the hallmark of embedding content that does not
    close on the combinatorial graph; (ii)~no three-term recurrence
    in $(n,l)$ exists for the integrals, in contrast to the Neumann
    expansion for $H_2$ in prolate spheroidal coordinates, where the
    separability of $1/r_{12}$ produces such a recurrence;
    (iii)~in the wavefunction-graph eigenbasis, 92--94\% of $V_{ee}$'s
    Frobenius mass is already diagonal
    ($\|[H_1, V_{ee}]\| / \|V_{ee}\|$ saturating at 5--6\% across
    $n_{\max} = 3, 4$), with the off-diagonal residual concentrated
    on the $(1s,1s)$ coalescence pair-state rather than distributed
    across angular channels.  The cusp is therefore localized at a
    single pair-state, and the wavefunction graph---which has no
    $1/r_{12}$ in its construction---nonetheless nearly diagonalizes
    it.  This is direct structural evidence that $1/r_{12}$ is an
    embedding exchange constant rather than an intrinsic graph object,
    and it explains why graph-native CI converges to a floor
    (Paper~13, Track~DI): the residual off-diagonal mass is precisely
    the embedding content the graph cannot absorb.

  \item \textbf{Algebraic (implicit) functions}: Functions of
    continuous parameters defined implicitly by $P(R,\mu) = 0$,
    where $P$ is the characteristic polynomial of the
    parameter-dependent Hamiltonian $H(R) = H_0 + R \cdot V_C$.
    The coefficients of $P$ lie in $\mathbb{Q}(\pi, \sqrt{2})$---a
    transcendental extension of $\mathbb{Q}$---but the $R$-dependence
    itself is algebraic: $\mu(R)$ is a root of a polynomial in $R$.
    The transcendental content is \emph{inherited} from the angular
    coupling matrix $V_C$ (which contains upstream exchange constants
    of embedding type), not generated by the parameterization.
    Example: $\mu(R)$ for the $\mathrm{SO}(6) \to \mathrm{SO}(4)$
    symmetry crossover on $S^5$ at Level~3.  At $l_{\max} = 0$,
    $\mu(R)$ is an explicit linear function of $R$; at higher
    $l_{\max}$, it is an algebraic function of degree $l_{\max}+1$,
    with branch points at eigenvalue crossings in the complex
    $R$-plane.  At Level~4, the split-region Legendre expansion of
    the nuclear charge function introduces piecewise-smooth
    $\rho$-dependence at the region boundary $\rho = 1/2$, breaking
    the linear pencil structure. The angular eigenvalues $\mu(\rho)$
    are piecewise-algebraic within each region but do not satisfy a
    single global characteristic polynomial. The product-space CI
    saturation at 92.5\% $D_e$ remains the key irreducibility result.

  \item \textbf{Composition constants}: The accuracy cost of
    factorizing a full $N$-electron wavefunction into composed
    coordinates.  Unlike the preceding types, which convert between
    discrete and continuous descriptions, composition constants
    convert between a single-geometry and a multi-geometry
    (factorized) description.  The factorization destroys inter-group
    antisymmetry---the exchange operator $P_{ij}$ mixes coordinate
    systems that composition separates---and the PK pseudopotential
    is the local approximation to this nonlocal effect.  Example:
    PK for LiH composed geometry (Level~3 core + Level~4 valence),
    which achieves 144$\times$ angular compression at the cost of
    ${\sim}5\%$ $R_{\rm eq}$ error at $l_{\max} = 2$.
\end{enumerate}

The dependency hierarchy is explicit: intrinsic constants depend
only on dimension; calibration constants depend on the graph;
embedding constants depend on the basis; flow functions depend
on the full parameterized geometry; composition constants depend
on the factorization topology (which groups share coordinates).

\paragraph{Odd-zeta content from first-order spectral operators (new tier).}
The Phase 4I Dirac-on-$S^3$ Tier 1 sprint (2026) identified a
structurally new kind of transcendental content in the framework: the
Apéry constant $\zeta_R(3)$, and more generally odd Riemann zetas
$\zeta_R(2k+1)$, arising natively from first-order spectral operators
(the Dirac operator) on curved space. The mechanism is the
Dirac-sector Dirichlet series at the packing exponent $s = d_{\max}$:
with the Camporesi--Higuchi spectrum $|\lambda_n| = n + 3/2$ and
degeneracy $g_n^{\text{Dirac}} = 2(n+1)(n+2)$ on unit $S^3$, the Fock
Dirichlet series evaluates to
\begin{equation}\label{eq:dirac_dirichlet_zeta3}
  D_{g_n^{\text{Dirac}}}^{\text{Fock}}(4)
  \;=\; \sum_{m \geq 1} \frac{2m(m+1)}{m^4}
  \;=\; 2\,\zeta_R(2) + 2\,\zeta_R(3)
  \;=\; \frac{\pi^2}{3} + 2\,\zeta_R(3),
\end{equation}
a closed-form symbolic identity (sympy-exact, see
\texttt{debug/dirac\_d3\_memo.md}). The two summands correspond to the
two homogeneity classes of $g_m^{\text{Dirac}} = 2m^2 + 2m$: the
weight-$m^2$ subchannel produces $2\zeta_R(2) = \pi^2/3$ (the familiar
scalar $F$ content of Paper~2), while the weight-$m$ subchannel produces
$2\zeta_R(3)$ (a channel structurally absent from the scalar Fock
degeneracy $g_n = n^2$, which has only one homogeneity class).

This odd-zeta content is categorically different from the
calibration-tier content cataloged above. Calibration constants
($\kappa$, $\pi^{d/2}$ Weyl prefactors, $F = \pi^2/6$ from the scalar
Fock Dirichlet series) carry \emph{even-zeta} or $\pi^{\text{even}}$
content arising from second-order Riemannian operators (the
Laplace--Beltrami and its conformal rescalings) through Jacobi-$\theta$
inversion in heat-kernel / Mellin analysis. The Phase 4D result
(Track~$\alpha$-H) established structurally that Hopf-fiber tensor
traces, heat-kernel asymptotics on round spheres, and spectral
determinants on odd-dimensional spheres all produce only integer (or
half-integer) powers of $\pi$ and rationals, not $\zeta_R(3)$. The
odd-zeta content $\zeta_R(3)$ first enters the framework through the
\emph{first-order} Dirac operator, via half-integer Hurwitz-zeta
identities in the Camporesi--Higuchi index, and is $\mathbb{Q}$-linearly
independent of $\zeta_R(2)$ by Apéry's theorem (irrationality 1978;
algebraic independence from $\pi$ remains open).

The appropriate taxonomic classification of $\zeta_R(3)$ is therefore a
\emph{new tier adjacent to calibration but categorically distinct}:

\begin{itemize}
  \item \textbf{Source:} first-order spectral operator on curved
    space (Dirac on $S^3$ in the present instance).
  \item \textbf{Mechanism:} Dirichlet series $\sum_n g_n \cdot
    (n+c)^{-s}$ at the packing exponent $s = d_{\max}$, where the
    shifted index $n + c$ ($c$ half-integer) and the mixed-homogeneity
    degeneracy together lift the scalar $\zeta_R(\text{even})$ result
    to $\zeta_R(\text{even}) + \zeta_R(\text{odd})$.
  \item \textbf{Dependency:} depends on the spectral sector (scalar
    vs.\ spinor) of the chosen operator, not on the discretization,
    basis, or factorization topology.
  \item \textbf{Algebraic independence:} $\zeta_R(3)$ is
    $\mathbb{Q}$-linearly independent of $\zeta_R(2)$ (Apéry), so no
    rational combination of Dirac-sector Dirichlet series reduces to
    a multiple of the scalar calibration content alone.
\end{itemize}

Within the $K = \pi(B + F - \Delta)$ combination rule of Paper~2 this
new tier is a \emph{byproduct rather than a summand}. $\zeta_R(3)$
appears in the Dirac Dirichlet series at $s = d_{\max}$ but is not used
by the formula for $\alpha^{-1}$; the formula uses only $B$ (scalar
Casimir trace), $F = \zeta_R(2)$ (scalar Fock Dirichlet), and
$\Delta^{-1} = g_3^{\text{Dirac}}$ (single-level spinor degeneracy). The
Phase 4I sprint confirmed that $\zeta_R(3)$ is structurally present
on $S^3$ without participating in $K$; see Paper~2 \S IV for the full
record. The taxonomic placement of $\zeta_R(3)$ as a distinct tier is
recorded here for use in future applications that may require
relativistic or spin-dependent spectral quantities.

\paragraph{Spinor-intrinsic content from first-order operators on a
spinor bundle (new subtier).}
The Dirac-on-$S^3$ Tier 2 sprint (2026, Track T5) identified a second
piece of structurally new transcendental content in the GeoVac
framework: the fine-structure constant squared~$\alpha^2$ and the
relativistic radial coefficient $\gamma := \sqrt{1 - (Z\alpha)^2}$,
arising jointly from first-order spectral operators acting on curved-space
\emph{spinor bundles}. The mechanism is the Breit--Pauli reduction of
the Dirac equation: the leading-order spin-orbit coupling
\begin{equation}\label{eq:bp_so}
  H_{\mathrm{SO}}
  \;=\; \tfrac{1}{2 m^2 c^2}\,\tfrac{1}{r}\,\tfrac{dV}{dr}\,\mathbf{L}\cdot\mathbf{S}
  \;=\; \tfrac{Z\alpha^2}{2}\,\tfrac{1}{r^3}\,\mathbf{L}\cdot\mathbf{S}
  \quad\text{(Coulomb),}
\end{equation}
when evaluated in the $(\kappa, m_j)$ basis using Szmytkowski's
reduction $\langle \mathbf{L}\cdot\mathbf{S}\rangle = -(\kappa+1)/2$ and
the closed-form hydrogenic $\langle 1/r^3\rangle_{n,l}
= Z^3/[n^3\, l(l+\tfrac{1}{2})(l+1)]$, produces matrix elements
\begin{equation}\label{eq:bp_so_matrix}
  \langle n,\kappa,m_j|\,H_{\mathrm{SO}}\,|n,\kappa,m_j\rangle
  \;=\; -\,\frac{Z^4\,\alpha^2\,(\kappa+1)}{4\,n^3\, l(l+\tfrac{1}{2})(l+1)},
\end{equation}
with $l = |\kappa + \tfrac{1}{2}| - \tfrac{1}{2}$ and the Kramers
cancellation $H_{\mathrm{SO}} = 0$ at $l=0$ (where $\kappa = -1$ forces
the numerator to zero before the formally divergent $\langle 1/r^3\rangle$
is evaluated). Symbolically, Eq.~(\ref{eq:bp_so_matrix}) is a sympy-exact
rational in $(Z, n, l)$ times $\alpha^2$, with zero $\pi$, zero $\zeta$,
and zero transcendentals beyond~$\alpha$ itself.

The full-Dirac radial corrections enter through
$\gamma := \sqrt{1 - (Z\alpha)^2}$, which appears in the Martínez-y-Romero
recursions for relativistic $\langle r^k\rangle_{\text{Dirac}}$ but is not
required at the Breit--Pauli leading order (Track T1 reserved $\gamma$
symbolically; Track T2 does not bind it). Algebraically, $\gamma$
satisfies the polynomial relation
\begin{equation}\label{eq:gamma_def}
  \gamma^2 + (Z\alpha)^2 \;=\; 1,
\end{equation}
so $\gamma \in \mathbb{Q}(\alpha^2)[\sqrt{\cdot}]$, i.e.\ it is
algebraic of degree~2 over $\mathbb{Q}(\alpha^2)$. This is the
relativistic analog of the non-relativistic radial matrix elements
already registered as algebraic in Sec.~\ref{sec:algebraic_registry}.

This spinor-intrinsic content is categorically different from the
four existing tiers (intrinsic, calibration, embedding, flow/composition)
and from the newly added odd-zeta tier. The distinctions:

\begin{itemize}
  \item \textbf{Source:} a first-order spectral operator (the Dirac
    operator or the Breit--Pauli reduction of it) acting on a
    \emph{spinor bundle} over curved space. Calibration constants and
    the odd-zeta tier both arise on scalar fields; the spinor-intrinsic
    content requires the bundle to carry spin.

  \item \textbf{Mechanism:} the $\alpha^2$ prefactor is not a basis
    regularization (as $\kappa$ is) and not a spectral sum (as
    $\zeta_R(2)$ and $\zeta_R(3)$ are); it is the coupling strength that
    multiplies the spin-orbit operator, arising from the Foldy--Wouthuysen /
    Breit--Pauli expansion of the relativistic kinetic energy. It is
    present in the Hamiltonian \emph{coefficients}, not in the basis or
    spectrum.

  \item \textbf{Dependency:} depends only on the existence of spin
    coupling in the operator. Independent of discretization, basis
    choice, and factorization topology. Any first-order operator that
    couples $\mathbf{L}$ and $\mathbf{S}$ will introduce $\alpha^2$;
    any full-Dirac radial recursion will introduce~$\gamma$.

  \item \textbf{Algebraic status:} $\alpha$ is transcendental (a
    dimensionless physical constant whose value must be supplied
    empirically), but $\alpha^2$ appears polynomially in every
    spinor-intrinsic coefficient to the order at which the Breit--Pauli
    expansion is truncated ($\alpha^2$ at leading order, $\alpha^4$ with
    Darwin + mass-velocity, etc.). $\gamma$ is algebraic of degree~2
    over $\mathbb{Q}(\alpha^2)$ via Eq.~(\ref{eq:gamma_def}). Together
    they generate the ring
    \[
      R_{\mathrm{sp}} \;:=\; \mathbb{Q}\bigl(\alpha^2\bigr)
        \bigl[\gamma\bigr]/\bigl(\gamma^2 + (Z\alpha)^2 - 1\bigr),
    \]
    in which every Tier-2 coefficient lies. A sympy-based certifier
    (``\texttt{geovac/spinor\_certificate.py}::verify\_spinor\_pi\_free})''
    mechanically checks that this ring is clean: no $\pi$, no $\zeta$,
    no $E_1$, no unregistered transcendental.
\end{itemize}

\paragraph{Updated structural partition.}
With the Tier 2 addition, the full Paper 18 partition of transcendental
content across bundle and operator-order now reads:

\begin{table}[h]
\caption{Operator-order $\times$ bundle classification of transcendental
  content in the GeoVac framework, after the Tier 1, Tier 2, and Tier 3
  Dirac-on-$S^3$ sprints.  The (2nd-order, spinor) cell is degenerate
  with the scalar calibration cell ($\pi^{\mathrm{even}}$ only).}
\label{tab:bundle_operator_tiers}
\begin{ruledtabular}
\begin{tabular}{llll}
  Tier & Example transcendentals & Operator order & Bundle \\
  \hline
  Intrinsic         & $\pi^{d/2}$ Weyl prefactors & --- & scalar \\
  Calibration       & $\kappa = -1/16$, $F = \pi^2/6$ (even $\zeta$)
                    & 2nd (Laplace--Beltrami)
                    & scalar on curved space \\
  Embedding         & $e^a E_1(a)$,  $1/r_{12}$
                    & ---
                    & projection from graph to $\mathbb{R}^3$ \\
  Flow / composition & piecewise $\mu(\rho,R)$, PK pseudopot.
                    & ---
                    & multi-geometry factorization \\
  Odd-zeta (Tier 1 new) & $\zeta_R(3), \zeta_R(5), \dots$
                    & 1st (Dirac)
                    & Dirac spectrum on curved space \\
  Spinor-intrinsic (Tier 2 new)
                    & $\alpha^2$, $\gamma = \sqrt{1-(Z\alpha)^2}$
                    & 1st (Dirac / Breit--Pauli)
                    & spinor bundle over curved space \\
  $D^2$ calibration (Tier 3 new)
                    & $\pi^{\text{even}}$ (degenerate w/ calibration)
                    & 2nd ($D^2 = \nabla^*\nabla + R/4$)
                    & spinor bundle over curved space \\
\end{tabular}
\end{ruledtabular}
\end{table}

The two Tier-1/Tier-2 additions are parallel but distinct: Tier 1
($\zeta$) enters through the \emph{spectrum} of the curved-space Dirac
operator (Camporesi--Higuchi eigenvalues $n+3/2$ and degeneracies
$2(n+1)(n+2)$ produce the Dirichlet-series identity
Eq.~(\ref{eq:dirac_dirichlet_zeta3})); Tier 2 ($\alpha^2$, $\gamma$)
enters through the \emph{coupling} of spin to orbital angular momentum
in the Hamiltonian, not through the basis or spectrum. Neither
introduces $\pi$ beyond what is already present in the calibration
tier; both are orthogonal to the embedding and flow tiers.

\paragraph{Why ``spinor-intrinsic'', not ``calibration''.}
Calibration constants live in the \emph{scalar} spectrum on curved
space (Paper~24 HO rigidity: the Bargmann--Segal edge weights are
$\pi$-free rationals; the calibration $\pi^2/6$ enters only through an
infinite-sum analysis of the scalar Fock Dirichlet series at $s = 4$).
Spinor-intrinsic $\alpha^2$ is a \emph{coupling} content --- it does
not appear in the basis construction or in the spectrum of the
operator; it appears in the Hamiltonian coefficients once spin is
coupled to orbital angular momentum. This distinction is taxonomic, not
cosmetic: it predicts that any future tier of the framework introducing
relativistic or QED corrections (e.g., $\alpha^4$ Darwin and
mass-velocity, $\alpha^3$ Lamb shift, Breit interaction
corrections $\alpha^2$ in $V_{ee}$) will populate this same subtier,
adding powers of $\alpha$ and algebraic combinations of $\gamma$, but
will \emph{not} create new subtiers and will \emph{not} bleed into
calibration or embedding content.

\paragraph{Certification and taxonomy wiring.}
The spinor-intrinsic ring $R_{\mathrm{sp}}$ is certified mechanically by
\texttt{verify\_spinor\_pi\_free} in
\texttt{geovac/spinor\_certificate.py}
(Track~T5, v2.11.0). The certifier walks the sympy expression tree of a
coefficient and rejects any appearance of $\pi$, $\zeta$, $\log$, $E_1$,
or any Symbol outside the registered set $\{Z, \alpha, \gamma\}$
(user-extensible for downstream indices). T3's symbolic spin-orbit block
\texttt{build\_so\_hamiltonian\_block(n\_max, Z\_sym, alpha\_sym)} passes the
certifier cleanly for all tested $n_{\max}$ up to~4; manual contamination
by $\pi \cdot \alpha^2 / 96$ is rejected with a clear diagnostic. See
the design memos \texttt{spin\_orbit\_design\_memo.md} §``Paper 18
classification'' and \texttt{spin\_ful\_composed\_design\_memo.md} §8 for
the per-block audit, and the sprint plan
\texttt{docs/dirac\_s3\_tier2\_sprint\_plan.md} Track~T5 for scope.

\paragraph{The squared Dirac operator and the completeness of the taxonomic grid.}
The Tier 3 Dirac-on-$S^3$ sprint (2026, Track T9) probed the last open
cell in the operator-order $\times$ bundle classification
(Table~\ref{tab:bundle_operator_tiers}): the squared Dirac operator
$D^2$ on the unit three-sphere, a second-order operator on spinor
sections.

By the Lichnerowicz formula~\cite{lichnerowicz1963}, $D^2 = \nabla^*\nabla + R/4$,
where $R = 6$ is the scalar curvature of unit $S^3$, so $R/4 = 3/2$ is a
rational constant shift.  The eigenvalues of $D^2$ are
$(n + 3/2)^2 = (2n+3)^2/4$ with degeneracies $g_n = 2(n+1)(n+2)$ (full
Dirac, both chiralities).  The spectral zeta function evaluates in closed
form at every integer $s$ via a partial-fraction identity: substituting
$m = n+1$ and using $m(m+1) = [(2m+1)^2 - 1]/4$,
\begin{equation}\label{eq:zeta_d2}
  \zeta_{D^2}(s)
  \;=\; 2^{2s-1}\bigl[\lambda(2s{-}2) - \lambda(2s)\bigr],
\end{equation}
where $\lambda(a) = (1 - 2^{-a})\,\zeta_R(a)$ is the Dirichlet lambda
function.  Since $\zeta_R(2k) = (-1)^{k+1}\,(2\pi)^{2k}\,B_{2k}/(2\cdot(2k)!)$
(Bernoulli formula), every term is a rational multiple of $\pi^{2k}$.
At every integer $s$, $\zeta_{D^2}(s)$ is therefore a two-term polynomial
in $\pi^2$ with rational coefficients:
$c_{2s-2}\,\pi^{2s-2} + c_{2s}\,\pi^{2s}$.  Explicitly:
$\zeta_{D^2}(1) = -\pi^2/4$ (analytic continuation),
$\zeta_{D^2}(2) = \pi^2 - \pi^4/12$,
$\zeta_{D^2}(3) = \pi^4/3 - \pi^6/30$,
$\zeta_{D^2}(4) = \pi^6(168 - 17\pi^2)/1260$.

\emph{No odd-zeta content appears at any~$s$.}  The structural mechanism
is algebraic: $D^2$ eigenvalues are perfect squares of odd half-integers,
so the spectral sums reduce to $\sum_k f(k)\,k^{-2s}$ over odd integers,
with exponent $2s$ (always even).  By the Bernoulli formula, such sums
are always $\mathbb{Q} \cdot \pi^{2s}$.  Squaring the operator maps the
first-order Dirichlet $\lambda(\text{odd})$ to $\lambda(\text{even})$,
structurally eliminating the odd-zeta content that Tier~1 (Track~D3)
found in the first-order Dirac spectral zeta via
Eq.~(\ref{eq:dirac_dirichlet_zeta3}).  The Hopf-equivariant per-charge
decomposition (half-integer charges for spinor modes) confirms this
universally across all Hopf sectors.

The (2nd-order, spinor-bundle) cell of Table~\ref{tab:bundle_operator_tiers}
is therefore \textbf{filled and degenerate with the scalar calibration
cell}: both produce $\pi^{\text{even}}$ content exclusively.  The
Lichnerowicz constant shift $R/4 = 3/2$ modifies rational coefficients
but does not introduce new transcendental content.  Operator order---first
vs.\ second---is the primary discriminant for transcendental class on
round~$S^3$.  Bundle type (scalar vs.\ spinor) modulates the
degeneracy lattice ($n^2$ vs.\ $2(n+1)(n+2)$), the eigenvalue lattice
(integer vs.\ half-integer), the Hopf charge lattice (integer vs.\
half-integer), and the rational coefficients in the $\pi^{2k}$
expansions, but it does \emph{not} change the transcendental class.

The taxonomic grid thus collapses to three effective tiers:
\begin{enumerate}
  \item \textbf{First-order operators:} produce odd-zeta content
    ($\zeta_R(3), \zeta_R(5), \ldots$) via sums over
    $(\text{odd})^{-s}$ at odd~$s$, plus coupling content ($\alpha^2$,
    $\gamma$) in the spinor sector.
  \item \textbf{Second-order operators:} produce $\pi^{\text{even}}$
    content exclusively, on both scalar and spinor bundles. This is the
    calibration tier of Paper~24.
  \item \textbf{Basis and composition tiers:} unchanged (embedding,
    flow, composition constants are operator-independent).
\end{enumerate}

\paragraph{Parallel to Paper 24.}
Paper~24 established that the Bargmann--Segal lattice is bit-exactly
$\pi$-free in rational arithmetic and that calibration $\pi$ is
structurally Coulomb-specific: the Coulomb/HO asymmetry arises because
second-order Riemannian operators (Coulomb via $-\Delta_{\text{LB}}$)
have nonlinear conformal projections that introduce $\pi$, whereas
first-order complex-analytic operators (the HO via the Bargmann Euler
operator on $H^2(S^5)$) have linear projections and no $\pi$.
Eq.~(\ref{eq:dirac_dirichlet_zeta3}) is the first-order, spinor analog
of that result on $S^3$: the Camporesi--Higuchi Dirac spectrum is
itself bit-exactly $\pi$-free in rational arithmetic (51/51 sympy tests
in \texttt{tests/test\_dirac\_s3.py}, v2.10.0), and $\pi^2$ enters
Eq.~(\ref{eq:dirac_dirichlet_zeta3}) only through the infinite-sum
analysis of the Dirichlet series at $s = 4$, not through the spectrum
itself. The $\zeta_R(3)$ content enters simultaneously and
inseparably. Cross-reference: Paper~24 \S IV (Coulomb/HO asymmetry)
for the second-order/scalar parallel; Paper~2 \S IV
(\S\,\texttt{zeta3\_aside}) for the $K$-combination-rule implications.

\paragraph{Motivic weight and operator order.}
The operator-order discriminant---second-order operators produce
$\pi^{\text{even}}$ content, first-order operators produce odd-zeta
content---has a natural interpretation in the theory of mixed Tate
motives over $\mathrm{Spec}(\mathbb{Z})$.  In this framework
(Goncharov~\cite{goncharov1999}, Zagier~\cite{zagier1994}), the Riemann
zeta value $\zeta_R(n)$ is a \emph{period} of the mixed Tate motive
$\mathfrak{M}(n) = \mathbb{Q}(-n)$ of weight~$n$.  Even zeta values
$\zeta_R(2k) = \text{rational} \times \pi^{2k}$ are periods of motives
that factor through the de~Rham cohomology of even-dimensional compact
manifolds (the Bernoulli formula).  Odd zeta values $\zeta_R(2k{+}1)$
are periods of genuinely independent motives: they do not reduce to
powers of~$\pi$, and $\zeta_R(3)$ in particular is irrational (Ap\'ery
1978) and conjectured algebraically independent from~$\pi$.

The GeoVac operator-order discriminant maps onto this motivic
partition:
\begin{center}
\begin{tabular}{llll}
  \hline
  Operator order & Motivic weight & Transcendental class & Example \\
  \hline
  2nd ($\Delta_{\text{LB}},\; D^2$)
    & even ($2k$) & $\zeta_R(2k) = \text{rat} \cdot \pi^{2k}$
    & $\kappa$, $\mathrm{Vol}(S^3)$, $a_2$ \\
  1st ($D$)
    & odd ($2k{+}1$) & $\zeta_R(2k{+}1) \perp \mathbb{Q}(\pi)$
    & D3: $\zeta_R(3)$ in Dirac Dirichlet \\
  \hline
\end{tabular}
\end{center}

At motivic weight~3, the appearance of $\zeta_R(3)$ in the GeoVac
framework (Eq.~(\ref{eq:dirac_dirichlet_zeta3}), Dirac Dirichlet series
on~$S^3$) belongs to a family of independent geometric sources that are
all periods of the same motivic object $\mathbb{Q}(-3)$:
\begin{enumerate}
  \item \emph{Dirac operator on $S^3$} (this work, Track~D3): the
    weight-$m$ subchannel of $g_m^{\text{Dirac}} = 2m^2 + 2m$ at
    $s = d_{\max} = 4$ produces $2\zeta_R(3)$.
  \item \emph{Hyperbolic 3-manifold volumes}: for arithmetic hyperbolic
    manifolds $\mathbb{H}^3/\Gamma$, the volume is a $\mathbb{Q}$-linear
    combination of values of the Bloch--Wigner dilogarithm, which are
    periods of $\mathbb{Q}(-3)$ (Thurston; Goncharov~\cite{goncharov1999}).
  \item \emph{Two-loop QED anomalous magnetic moment}: $\zeta_R(3)$
    enters the electron $g{-}2$ at order $\alpha^2$ through irreducible
    nested harmonic sums in Feynman parameter integrals
    (Petermann 1957; Sommerfield 1957).
  \item \emph{Chern--Simons theory on $S^3$ at two loops}: the
    perturbative Chern--Simons invariant acquires $\zeta_R(3)$ at
    the two-loop level.
  \item \emph{Algebraic K-theory}: $\zeta_R(3)$ is a period of the
    Borel regulator on $K_5(\mathbb{Z})$, connecting number theory to
    the motivic cohomology of $\mathrm{Spec}(\mathbb{Z})$.
\end{enumerate}
The common thread is the three-sphere: items 1, 2, and~4 involve $S^3$
or its quotients directly; item~3 involves products of massless
propagators on~$S^3$ (after Wick rotation); item~5 is the number-theoretic
shadow of the $S^3$ geometry via the Borel regulator.

Zagier's dimension conjecture~\cite{zagier1994} predicts that the
$\mathbb{Q}$-vector space of weight-$n$ multiple zeta values has
dimension $d_n$ satisfying $d_0 = 1$, $d_1 = 0$, $d_2 = 1$,
$d_n = d_{n-2} + d_{n-3}$.  At weight~3 this gives $d_3 = 1$:
$\zeta_R(3)$ is the unique independent transcendental at that weight.
This is consistent with the GeoVac finding that $\zeta_R(3)$ appears
exactly once (in the Dirac sector, Tier~1) and is absent from the
scalar sector at every order.

The connection to QED loop order is structurally predicted by the
motivic correspondence.  The T9 theorem
(Eq.~(\ref{eq:zeta_d2}), $\zeta_{D^2}(s)$ involves only
$\pi^{\text{even}}$) guarantees that $\zeta_R(3)$ cannot enter through
any single trace of a squared propagator $D^{-2s}$---i.e., at one loop.
It enters at two loops, where irreducible products of first-order
propagators $D^{-1} \otimes D^{-1}$ produce odd motivic weight.  The
operator-order discriminant thus determines the \emph{loop order} at
which new transcendental content first appears: one loop accesses only
even-weight motives ($\pi^{2k}$); two loops access odd-weight motives
($\zeta_R(3)$, $\zeta_R(5)$, \ldots).

This motivic interpretation refines the compactness thesis of
Sec.~\ref{sec:intro}: exchange constants are not only projections from
compact to non-compact geometry---they are \emph{periods} of specific
mixed Tate motives, and the operator order on the compact space
determines the motivic weight of the resulting transcendental content.

\paragraph{Irreducibility of the Level~4 piecewise structure.}
The piecewise-smooth $\rho$-dependence of $\mu(\rho,R)$ at Level~4
is not a coordinate artifact: it cannot be resolved by geometric
elevation to any higher-dimensional space (Track~Z).  Three approaches
were ruled out: (i)~algebraic blow-up of $|s - \rho|$ produces smooth
sheets but with $\rho$-dependent integration limits ($\alpha = \arccos\rho$,
$\arcsin\rho$), making matrix elements transcendental in~$\rho$;
(ii)~Lie algebra elevation fails because $\min(s,\rho)/\max(s,\rho)$
is $C^0$ but not~$C^1$, incompatible with matrix elements of any
finite-dimensional Lie group representation (which are real-analytic);
(iii)~$S^3 \times S^3$ product space fails because two-center nuclear
attraction cannot be smooth on any single compact manifold.
The piecewise structure thus belongs to the embedding category
in the taxonomy: it is the irreducible geometric content of placing
two nuclei at finite separation in hyperspherical coordinates, analogous
to the $1/r_{12}$ codimension obstruction (Track~W) but affecting the
nuclear rather than the electronic coupling.

\paragraph{Even-odd staircase and channel convergence rate.}
An extended convergence study at Level~4 (Paper~15, Track~AA)
reveals a pronounced even-odd staircase in $D_e$ vs $l_{\max}$:
even increments produce large gains ($+42$, $+6.8$, $+1.7$~pp)
while odd increments produce near-zero gains ($+0.66$, $+0.08$~pp).
The root cause (Track~AC) connects directly to the exchange constant
structure: the nuclear coupling matrix
$\langle l_1' l_2' | V_{\rm nuc} | l_1 l_2 \rangle$ is diagonal
in one electron's angular momentum index (Kronecker delta), so
direct coupling from the ground channel $(0,0)$ requires a
$\Delta l = 2$ transition enforced by the gerade constraint
$l_1 + l_2 = \text{even}$.  Odd-$l_{\max}$ channels couple only
indirectly through the electron-electron multipole expansion.

This selection rule constrains the \emph{convergence rate} of the
partial-wave expansion: the nuclear coupling---a piecewise exchange
constant in the taxonomy---couples to the ground state only through
quadrupolar ($\Delta l = 2$) transitions, producing an algebraic
convergence rate $\sim (l + 1/2)^{-2}$ with an even-odd modulation.
The $\sigma$ and $\pi$ channel sectors are completely decoupled
(both nuclear and $e$--$e$ coupling conserve $m$), with $\pi$
contributing a constant ${\sim}6.6$~pp additive offset.

\paragraph{PK pseudopotential as composition exchange constant.}
The Phillips--Kleinman pseudopotential in the composed geometry
(Paper~17) enforces core-valence orthogonality, which is a
Pauli exclusion effect---distinct from the Coulomb exchange
constants cataloged above.  Without PK, no equilibrium
geometry exists (the PES is monotonically attractive), establishing
PK as essential for creating the bonding minimum.  Ab initio
derivation of the $R$-dependent PK scaling parameter $R_{\rm ref}$
from core properties fails due to a fundamental scale mismatch:
all core-derived candidates ($0.27$--$0.82$~bohr) are too small
by a factor of $4$--$11$.

A complete survey of all solver $\times$ PK combinations
(Paper~17, Table~VI) confirms that PK is not an improvable
approximation.  The 2D variational solver eliminates ${\sim}75\%$
of the $l_{\max}$ drift (from the adiabatic approximation), but
$l$-dependent PK has zero additional effect on the 2D solver.
The residual drift ($+0.15$~bohr/$l_{\max}$) is intrinsic to
the composed-geometry factorization.

Three graph-native alternatives to PK were investigated and all
failed: (i)~node exclusion between the core $S^3$ and valence
Level~4 graphs has no coordinate correspondence;
(ii)~the fiber bundle Berry connection is intra-group, and the
exchange operator $P_{ij}$ mixes coordinate systems;
(iii)~spectral orthogonality (function-space) does not imply
coordinate-space exclusion.  The master obstruction is that
antisymmetry requires a shared coordinate system; composition
factorizes $\Psi_{\rm total}$ into incompatible coordinates.

PK therefore defines a sixth exchange type in the taxonomy:
\textbf{composition constants} convert between a single-geometry
(full $N$-electron) and a composed-geometry (factorized) description.
Unlike the other exchange constants, which convert between discrete
and continuous descriptions, composition constants measure the
accuracy cost of the factorization itself.  The ``exchange rate''
is quantified by the 144$\times$ angular compression that composed
LiH achieves over a full 4-electron solver on $S^{11}$
(Paper~17, Sec.~V): the ${\sim}5\%$ $R_{\rm eq}$ error at
$l_{\max} = 2$ is the irreducible cost of this compression.

This classification is validated by the full 4-electron Level~4N
solver (Paper~17, Sec.~VIII): exact $S_4$ antisymmetry produces
molecular equilibrium without PK at $l_{\max} \geq 2$, confirming
that PK approximates---rather than creates---the physics of Pauli
exclusion.  The composition exchange constant is the cost of
replacing exact inter-group antisymmetry with a pseudopotential
in the factorized coordinate system.


%% ========================================================================
%% V. THE FINE STRUCTURE CONSTANT CONNECTION
%% ========================================================================

\section{Connection to the Fine Structure Constant}
\label{sec:alpha}

Paper~2~\cite{loutey_paper2} observes that the fine structure
constant $\alpha$ satisfies the cubic $\alpha^3 - K\alpha + 1 = 0$,
where the coefficient
\begin{equation}
  K = \pi\!\left(42 + \frac{\pi^2}{6} - \frac{1}{40}\right)
  = 137.036064
  \label{eq:K}
\end{equation}
is constructed from spectral invariants of the Hopf fibration
$S^1 \to S^3 \to S^2$.  The combination rule
$K = \pi(B + F - \Delta)$ was identified as conjectural---an
empirical formula with structural support but no first-principles
derivation.

We observe that this combination rule has precisely the structure
of a spectral-geometric exchange:
\begin{itemize}
  \item $B = 42$ is discrete representation-theoretic data:
    the degeneracy-weighted $\mathrm{SO}(3)$ Casimir trace over
    the first three shells of $S^3$.
  \item $F = \zeta(2) = \pi^2/6$ is a continuous spectral
    invariant: the spectral zeta function of $S^1$ at $s = 1$.
  \item $\Delta = 1/40$ is a boundary correction mixing discrete
    eigenvalue data ($|\lambda_3| = 8$) with the cumulative state
    count ($N(2) = 5$).
  \item The overall factor $\pi$ is $\omega_2 = \pi$, the 2D
    ball volume, which serves as the numerator of the $S^2$ Weyl
    density and as the packing-plane prefactor
    $\sigma_0 = \pi d_0^2 / 2$ in Paper~0 (the base of the Hopf
    fibration, the photon propagation sphere).
\end{itemize}

The structure of $K$ is thus: the 2D ball volume $\omega_2 = \pi$
(the numerator of the $S^2$ Weyl density, and the same $\pi$ that
fixes Paper~0's packing-plane prefactor $\sigma_0 = \pi d_0^2 / 2$)
multiplying a sum of the $S^2$ Casimir trace ($B$, discrete),
the $S^1$ spectral zeta value ($F$, transcendental), and an
$S^3$ boundary term ($\Delta$, mixed).  This is a Weyl-type
formula that converts between the discrete bundle
(representation labels on $S^3$) and the continuous coupling
(the electromagnetic interaction strength measured in
$\mathbb{R}^3$).

If this identification is correct, then the combination rule
$K = \pi(B + F - \Delta)$ is not arbitrary numerology but an
instance of the same spectral-geometric exchange that produces
$\kappa = -1/16$ at Level~1 and $e^a E_1(a)$ at Level~2.
The exchange constant for the Hopf bundle is $K$ itself---the
``cost'' of projecting the $S^3$ electron state space through
the $S^1$ gauge fiber onto the $S^2$ photon base, measured
in units of coupling strength.

\begin{conjecture}
The combination rule $K = \pi(B + F - \Delta)$ is a
spectral-geometric exchange constant for the Hopf fibration,
analogous to the Weyl--Selberg constants for compact manifolds.
The specific form---$\pi$ multiplying a sum of Casimir data
and zeta values---should be derivable from the spectral geometry
of $S^1 \to S^3 \to S^2$, potentially via the heat kernel
expansion or the analytic torsion of the Hopf bundle.
\end{conjecture}

Two additional observations support this identification:

\begin{observation}[Structural match]
The combination rule $K = \pi(B + F - \Delta)$ mixes discrete
data ($B = 42$, from representation theory) with continuous
spectral data ($F = \zeta(2) = \pi^2/6$, from the $S^1$
spectral zeta function), multiplied by an overall geometric
factor ($\omega_2 = \pi$, the 2D ball volume that enters as the
numerator of the $S^2$ Weyl density and as the packing-plane
prefactor $\sigma_0 = \pi d_0^2 / 2$ in Paper~0, with $S^2$ the
Hopf base).  This is precisely the pattern identified throughout this
paper: a transcendental conversion factor mediating between
discrete and continuous descriptions of the same geometry.
The structural match is evidence---not proof---that the
combination rule is an instance of the exchange constant pattern
rather than numerology.
\end{observation}

\begin{observation}[$\alpha$ is likely transcendental]
The fine structure constant $\alpha$ is a root of the cubic
$\alpha^3 - K\alpha + 1 = 0$, where $K$ involves $\pi$ and
$\zeta(2) = \pi^2/6$.  Roots of polynomials with transcendental
coefficients \textit{can} be algebraic (e.g., $x^2 - \pi x +
\pi(\pi - 1) = 0$ has the algebraic root $x = 1$), but this
requires special cancellations that are generically absent.
Since $K$ is a non-trivial combination of $\pi$ and $\pi^2$
with rational coefficients, it is overwhelmingly likely
(though not proven) that $\alpha$ is transcendental---consistent
with its role as a spectral-geometric exchange constant.
\end{observation}

\subsection{Three-tier structure of the combination rule}
\label{sec:alpha_three_tier}

An ongoing computational investigation into the origin of the three
components of $K = \pi(B + F - \Delta)$---using exact-rational
symbolic computation and high-precision arithmetic---has pinned down
$B$ and $F$ as genuinely different objects on the Fock $(n,l)$
lattice, and has eliminated a broad family of candidate mechanisms
for $\Delta$. We observe that the three terms appear to occupy
three categorically distinct tiers of the exchange-constant
taxonomy rather than specializing a single common generator. We
summarize the four structural findings below; none of them alters
the conjectural status of Paper~2's cubic identity.

\paragraph{1. A $\kappa$--$B$ identity from the Fock weight.}
The Fock projection of the hydrogenic basis onto $S^3$ carries a
residual momentum-space weight
$w(p) = (p^2 + p_0^2)^{-2}$
after the overall $(2p_0)^{5/2}$ prefactor of the Fock map is
stripped.  Computation reveals that the diagonal matrix elements
$\langle n,l|w|n,l\rangle_{p_0 = 1}$ for $n \le 3$ are exact
rationals whose denominators divide $32 = 2\,d_{\max}^2$:
\begin{equation}
  \{\langle n,l|w|n,l\rangle_{p_0 = 1}\}_{(n,l)}
  = \Bigl\{ \tfrac{7}{16},\, \tfrac{5}{8},\, \tfrac{3}{8},\,
            \tfrac{5}{8},\, \tfrac{17}{32},\, \tfrac{11}{32} \Bigr\}
  \label{eq:fock_weight_values}
\end{equation}
for $(n,l) \in \{(1,0),(2,0),(2,1),(3,0),(3,1),(3,2)\}$.  Weighted
by the Hopf base factor $(2l+1)\,l(l+1)$ and summed,
\begin{equation}
  \sum_{n=1}^{3}\sum_{l=0}^{n-1} (2l+1)\,l(l+1)\,
    \langle n,l|w|n,l\rangle_{p_0 = 1}
  = \frac{63}{4}
  = 6\,B\,|\kappa|,
  \label{eq:kappa_B_identity}
\end{equation}
with $B = 42$ the Casimir trace of Paper~2 and $\kappa = -1/16$
the calibration constant of Sec.~\ref{sec:kappa_derivation}.
Equivalently, averaging over the six $(n,l)$ cells gives
$B\,|\kappa| = 21/8$.  This is the first exact rational algebraic
link in the GeoVac framework between the calibration constant and
the Hopf base invariant: $|\kappa| = 1/d_{\max}^2$ appears as the
universal denominator factor of the Fock-weight matrix elements,
while $B$ appears in the Hopf-weighted sum.  The identity
(\ref{eq:kappa_B_identity}) is verified by exact-rational symbolic
computation with no fitted parameters.

\paragraph{2. A Dirichlet identity for $F$.}
Define the Dirichlet series of the $S^3$ Fock shell degeneracy
$g_n = n^2$,
\begin{equation}
  D_{n^2}(s) \;:=\; \sum_{n=1}^{\infty} g_n\, n^{-s}
             \;=\; \sum_{n=1}^{\infty} n^{-(s-2)}
             \;=\; \zeta_R(s - 2),
  \label{eq:Dn2_def}
\end{equation}
evaluated at the packing exponent $s = d_{\max} = 4$ of Paper~0:
\begin{equation}
  D_{n^2}(d_{\max}) = \zeta_R(2) = \frac{\pi^2}{6} = F.
  \label{eq:F_dirichlet}
\end{equation}
The exponent $s = 4$ admits three independent natural readings:
$s = d_{\max}$ from Paper~0's packing axiom; $s = \dim(\mathbb{R}^4)$,
the ambient Euclidean space into which $S^3$ is stereographically
projected in Paper~7; and $s = 2 N_{\mathrm{init}}$ from the same
packing axiom.  The weight $n^2$ is the degeneracy of the $n$-th
$S^3$ shell in Fock's projection.  To our knowledge
(\ref{eq:F_dirichlet}) is the first identification of $F$ from a
graph-intrinsic quantity in the GeoVac framework---that is, $F$ is
produced here without invoking an embedding metric or a
sphere-spectral construction.  No other natural Dirichlet series
on the shell lattice (per-shell Casimir, state count, Laplacian
eigenvalue, $l$-max weight) produces an exact hit at any integer
$s$ in the range $[3, 12]$.

\paragraph{3. Three-tier classification.}
Combining the $\kappa$--$B$ identity (\ref{eq:kappa_B_identity})
with the Dirichlet identity (\ref{eq:F_dirichlet}) and the
Paper~2 canonical form $1/\Delta = |\lambda_{n_{\max}}|\, N(n_{\max}{-}1)$,
the three components of $K = \pi(B + F - \Delta)$ have
categorically different origins:

\begin{center}
\small
\renewcommand{\arraystretch}{1.2}
\begin{tabular}{@{}lccl@{}}
\hline
Component & Value & Type & Origin \\
\hline
$B$      & $42$        & rational, finite sum       & Casimir trace
   truncated at $n_{\max}=3$ \\
$F$      & $\pi^2/6$   & transcendental, infinite   & Fock-degeneracy
   Dirichlet at $s = d_{\max}$ \\
         &             & series                     & \\
$\Delta$ & $1/40$      & rational, boundary product & $|\lambda_{n_{\max}}|
   \cdot N(n_{\max}{-}1)$ \\
\hline
\end{tabular}
\end{center}

\noindent The three objects do not arise from a single common
generator.  $B$ is a finite sum over the first three shells;
$F$ is an infinite Dirichlet series with no meaningful
finite truncation at $n_{\max} = 3$ (the tail $\pi^2/6 - 49/36
\approx 0.284$ is not $\Delta$, nor any rational multiple of it);
and $\Delta$ is a boundary product at the truncation edge.
The combination rule $K = \pi(B + F - \Delta)$ therefore assembles
objects from three distinct tiers of the exchange-constant
taxonomy rather than specializing a single common Dirichlet
series, spectral construction, or arithmetic generator.

\paragraph{4. Negative results on candidate mechanisms.}
The three-tier reading is supported by the systematic elimination
of alternative unifying mechanisms.  Computation (exact rational
symbolic arithmetic and \texttt{mpmath} at 50 decimal places)
has ruled out, within the tolerance indicated:
(i) the Hopf-twist spectral comparison $S^3$ versus $S^1 \times S^2$
(cleanest near-miss $\sim 6.8 \times 10^{-3}$);
(ii) higher Casimir traces on $S^3$ at all orders through
$n_{\max} = 3$;
(iii) Hopf graph quotient Laplacian invariants for
$S^3 \to S^2$ over the $m$-fiber;
(iv) the discrete $S^1$ fiber spectral zeta (path-graph Laplacian,
in four weighting schemes);
(v) the continuous $S^1$ fiber via Mellin and heat-kernel expansions
(the Jacobi inversion $\theta_{S^1}(t) \sim \sqrt{\pi/t}$ forces every
order to be linear in $\pi$: the $\sqrt{\pi}\cdot\sqrt{\pi}\cdot\mathbb{Q}
= \pi\cdot\mathbb{Q}$ structure collapses any expected $\zeta_R(2)$
contribution to $\pi$ times a rational at every order);
(vi) $S^5$ spectral geometry, including zetas, Seeley--DeWitt
coefficients, and functional determinants (all rational, or
reducible to $\zeta_R(\mathrm{odd})$ and $\log 2$);
and (vii) a $\zeta$-combination scan for $\Delta$ over $3{,}228$
candidates built from Riemann and Hurwitz zetas, Bernoulli numbers,
and Stieltjes constants, with zero strict hits within $10^{-25}$
of $1/40$.

The structural reason for these negative results is that
round-sphere Laplace--Beltrami operators produce (i) integer powers
of $\pi$ in their volumes, (ii) rational Seeley--DeWitt
coefficients, and (iii) $\zeta_R(\mathrm{odd}) + \log 2$ combinations
in their spectral determinants---but never $\pi^2 = 2\zeta_R(2)$
additively next to a rational.  This is consistent with the
$\pi$-free certificate of the Bargmann--Segal lattice
(Paper~24), which shows that the
analogous lattice for the 3D harmonic oscillator on the Hardy
space of $S^5$ is bit-exactly rational in exact arithmetic; and the
same rigidity holds for $S^5$ Laplace--Beltrami spectral invariants
generally.

\paragraph{Reframed open question.}
The derivation question ``where does each piece of $K$ come from?''
has been partially answered: $B$ from a finite Casimir truncation
(with the $\kappa$--$B$ link of Eq.~\ref{eq:kappa_B_identity}), $F$
from an infinite Dirichlet series at the packing exponent
(Eq.~\ref{eq:F_dirichlet}), and $\Delta$ from a boundary product at
the truncation edge (irreducible to $\zeta$-content within the
candidate set tested).  The remaining structural question is no
longer a derivation question about any single component: it is a
question about why the \emph{additive} combination
$B + F - \Delta$, multiplied by the overall factor $\omega_2 = \pi$,
equals $\alpha^{-1}$ to $8.8 \times 10^{-8}$.  That is a question
about the conspiracy of three categorically different objects---a
finite sum, an infinite Dirichlet series, and a boundary
product---rather than the specialization of a single generator.
Paper~2's identification remains conjectural.


%% ========================================================================
%% VI. OBSERVABLE CLASSIFICATION
%% ========================================================================

\section{Observable Classification by Transcendental Content}
\label{sec:observable_classification}

The exchange constant taxonomy (Sec.~\ref{sec:taxonomy}) classifies the
\emph{constants themselves} by their mathematical origin.  A complementary
classification addresses the \emph{observables}: given a physical quantity
to be computed, what class of transcendental content does it require?

The conformal equivalence between the graph and continuous descriptions
(Paper~7) implies a natural partition of observables into three classes,
determined by what type of projection is needed to evaluate them.

\begin{enumerate}
  \item \textbf{Class~S (spectral observables):} Energy eigenvalues,
    degeneracies, selection rules, and transition amplitudes between
    eigenstates.  These require only the graph eigenvalue spectrum
    and the calibration constant $\kappa$.  No transcendental content
    beyond $\kappa = -1/16$ (rational) is needed.

  \item \textbf{Class~P (single-particle spatial observables):}
    Electron density $\rho(\mathbf{r})$, expectation values
    $\langle r^k \rangle$, momentum distributions, and any quantity
    requiring the wavefunction evaluated at a point in $\mathbb{R}^3$.
    These require the conformal projection $S^3 \to \mathbb{R}^3$,
    which introduces $\pi$ through the volume element
    $d\Omega_{S^3} = 2\pi^2$ and the spherical harmonic
    normalization $\int |Y_{lm}|^2 d\Omega = 1$ (with
    $d\Omega$ carrying $4\pi$).

  \item \textbf{Class~C (multi-particle correlated observables):}
    Correlation energies, adiabatic potential energy surfaces,
    multi-electron expectation values, and any quantity requiring
    the inter-particle coordinate.  These require the higher exchange
    constants: $e^a E_1(a)$ at Level~2 for $m \neq 0$ states
    (Sec.~\ref{sec:stieltjes}), $\mu(R)$ at Level~3
    (Sec.~\ref{sec:mu}), $\mu(\rho,R)$ at Level~4
    (Sec.~\ref{sec:mu_level4}), and the PK composition constant
    at Level~5.
\end{enumerate}

\noindent This classification is stated as a structural claim:

\medskip
\noindent\textbf{Claim~4.}
\textit{The class of transcendental content required to evaluate a
quantum-mechanical observable is determined by the type of projection
linking the graph description to the coordinate system in which the
observable is defined.  Spectral observables (Class~S) require no
transcendental content.  Single-particle spatial observables (Class~P)
require the Weyl constants ($\pi^{d/2}$).  Multi-particle correlated
observables (Class~C) require the embedding, flow, or composition
exchange constants of Sec.~\ref{sec:taxonomy}.}

\medskip

Claim~4 is falsifiable: for any observable, one can determine
its class, identify the required exchange constants, and verify
that no additional transcendental content enters the computation.
The following worked examples illustrate each class.

\subsection{Example 1: Hydrogen energy levels (Class~S)}

The hydrogen energy spectrum $E_n = -Z^2/(2n^2)$ is a spectral
observable.  The graph Laplacian on $S^3$ has eigenvalues
$\lambda_n = -(n^2 - 1)$, which are integers for all $n$
(Paper~7, Eq.~(12); Paper~0, Sec.~II).  The degeneracies
$g_n = n^2$ are integers.  The Fock conformal mapping (Paper~7,
Sec.~IV) converts these integer eigenvalues to the Rydberg
spectrum via the energy-shell constraint $p_0^2 = -2E$
and the stereographic focal length:
\begin{equation}
  E_n = -\frac{Z^2}{2n^2}\;\text{Ha},
  \label{eq:hydrogen_class_s}
\end{equation}
with the calibration constant $\kappa = -1/16$ (rational)
mediating between graph eigenvalues and physical energies
(Sec.~\ref{sec:kappa_derivation}).  Every quantity in this
derivation---eigenvalues, degeneracies, $\kappa$---is an
integer or rational number.  No transcendental content enters.

Selection rules provide a second example.  The dipole selection
rule $\Delta l = \pm 1$ follows from the triangle inequality
on the Gaunt integral
$\int Y_{l_1 m_1} Y_{1q} Y_{l_2 m_2}\,d\Omega$, which
reduces to a product of Wigner $3j$ symbols---rational numbers
computed from factorials of integers~\cite{loutey_paper0}.
No transcendental content enters.

\subsection{Example 2: Electron density (Class~P)}

The electron density $\rho(\mathbf{r}) = |\psi(\mathbf{r})|^2$
requires evaluating the wavefunction in position space.  On the
graph, the state $|n, l, m\rangle$ is a node with an integer
label; in $\mathbb{R}^3$, the corresponding wavefunction is
$\psi_{nlm}(\mathbf{r}) = R_{nl}(r)\,Y_{lm}(\theta,\phi)$.
The conformal projection $S^3 \to \mathbb{R}^3$ (Paper~7,
Sec.~III) introduces the conformal factor
$\Omega = 2p_0/(p^2 + p_0^2)$, which mediates between the
unit-sphere measure $d\Omega_{S^3}$ and the flat-space volume
element $d^3\mathbf{r}$.  The spherical harmonic normalization
$\int |Y_{lm}|^2\,d\Omega = 1$ introduces $\pi$ through the
solid angle $\oint d\Omega = 4\pi$.  The radial normalization
$\int_0^\infty |R_{nl}|^2\,r^2\,dr = 1$ introduces no additional
transcendentals for hydrogen (the integrals are rational multiples
of $a_0^{-3}$ times exponentials, evaluated in closed form).

The density is thus a Class~P observable: it requires the $\pi$
that enters through the projection from the dimensionless $S^3$
topology to the dimensionful $\mathbb{R}^3$ coordinate system.

\subsection{Example 3: Helium ground state energy (Class~C)}

The helium ground state energy $E_0 = -2.9037$~Ha~\cite{drake2006}
is a correlated observable.  The Level~3 hyperspherical solver
(Paper~13) computes it via the adiabatic potential $U(R)$, which
depends on the angular eigenvalue $\mu(R)$.  The function $\mu(R)$
satisfies the characteristic polynomial $P(R,\mu) = 0$
(Sec.~\ref{sec:algebraic_curve}), with coefficients in
$\mathbb{Q}(\pi, \sqrt{2})$ inherited from the angular coupling
matrix~\cite{loutey_track_p1}.  The energy is obtained by
integrating the hyperradial Schr\"odinger equation with
$U(R)$ as the effective potential---a definite integral over
the algebraic curve, producing a transcendental number.

On the $S^3 \times S^3$ product space (Level~3 FCI), the same
energy can in principle be computed without $\mu(R)$: the
Hamiltonian matrix has entries given by Gaunt integrals (rational)
and the Slater $F^0$ direct Coulomb integrals (algebraic on $S^3$,
Paper~7 Sec.~VI).  The FCI result (0.35\% error) demonstrates that
the correlated energy \emph{can} be approximated within the
algebraic framework, but the $\mu(R)$ parameterization is needed
to achieve sub-0.1\% accuracy.  The transition from Class~S to
Class~C reflects the physical content: electron-electron
correlation couples the two $S^3$ factors, and the coupling
strength varies with the hyperradial coordinate $R$---a spatial
degree of freedom that exists only in the projected description.

\paragraph{Three-layer decomposition (v2.6.0).}
The algebraic Casimir CI (Track~DI Sprint~2) makes the exchange
constant structure of the He ground state fully explicit.  At each
$n_{\max}$, the FCI matrix $H(k) = Bk + Ck^2$ has rational entries
as a function of the orbital exponent~$k$, so the variationally
optimal energy $E^*$ is an algebraic number---a root of a polynomial
with rational coefficients.  At $n_{\max} = 1$: $k^* = 27/16$,
$E^* = -729/256$~Ha (exact rationals).  This reveals a three-layer
structure:
\begin{center}
\begin{tabular}{lll}
\hline
Layer & Content & Exchange type \\
\hline
Rational & Slater integrals, Gaunt coefficients & None (graph-intrinsic) \\
Algebraic & Optimal exponents at finite $n_{\max}$ & Calibration \\
Transcendental & e-e cusp ($\alpha = \pi/4$ on $S^5$) & Embedding \\
\hline
\end{tabular}
\end{center}
The Fock energy-shell self-consistency $k^2 = -2E$, which is exact
at Level~1 (one electron, one scale), over-constrains the
two-electron problem: self-consistent~$k$ gives 5--13\% error vs.\
1.5--1.9\% for variational~$k$.  This is the first instance where a
single calibration constant is insufficient---the transition from
calibration to embedding exchange constants corresponds to the
physical content of electron-electron correlation.

\subsection{Boundary cases}

Two boundary cases clarify the classification:

\begin{itemize}
  \item \textbf{Ionization energies} are spectral (Class~S) when
    computed as eigenvalue differences $\Delta E = E_n - E_m$,
    which involve only $\kappa$ and integer quantum numbers.
    However, photoionization \emph{cross sections} require the
    dipole matrix element $\langle \psi_f | \mathbf{r} | \psi_i
    \rangle$, which is a spatial observable (Class~P) because
    it evaluates the wavefunction at positions in $\mathbb{R}^3$.

  \item \textbf{Two-electron energies} computed via FCI on
    $S^3 \times S^3$ are formally Class~S (graph eigenvalues
    of the product-space Hamiltonian), but achieve limited
    accuracy due to the product-basis completeness deficiency
    (Sec.~\ref{sec:mu_level4}).  High-accuracy two-electron
    energies require the Class~C exchange constants.
    The classification thus tracks both the type of observable
    and the accuracy target.
\end{itemize}

The classification does not assert that one class is more
fundamental than another.  The conformal equivalence
(Paper~7) ensures that all three descriptions---graph, projected
single-particle, correlated multi-particle---are mathematically
equivalent.  The classification identifies which exchange
constants mediate between them.


%% ========================================================================
%% VII. DISCUSSION
%% ========================================================================

\section{Discussion}
\label{sec:discussion}

\subsection{Why the constants are transcendental}

The fundamental reason that spectral-geometric exchange constants
are transcendental is that they convert between two
incommensurable descriptions: a discrete count (eigenvalues,
representation labels, graph vertices) and a continuous measure
(volume, curvature, geodesic length).  The Weyl constants involve
$\pi^{d/2}$ (Sec.~\ref{sec:pattern}), and $\pi$ is
transcendental by Lindemann's theorem (1882: $e^{i\pi} = -1$
implies $\pi$ is transcendental, since $e^\alpha$ is
transcendental for nonzero algebraic $\alpha$).  Since $\pi$
enters every volume-to-eigenvalue conversion, transcendence is
a structural feature of the discrete-to-continuous exchange, not
an accident of particular
manifolds~\cite{berger2003}.

However, not all \emph{functions} that appear in the GeoVac hierarchy
are transcendental.  The adiabatic eigenvalues $\mu(R)$ satisfy a
polynomial equation $P(R,\mu) = 0$ (Sec.~\ref{sec:algebraic_curve}),
making them algebraic functions of $R$.  The transcendental content
in $\mu(R)$ is entirely inherited from the coefficients of $P$, which
involve $\pi$ and $\sqrt{2}$ from the angular coupling integrals.
The $R$-dependence itself is algebraic.  This distinction matters:
the truly irreducible transcendentals are the \emph{constants} that
enter the coupling matrix, not the parameterized functions built
from them.

\subsection{Why the constants vanish on the graph}

When the computation remains entirely on the discrete graph---no
projection, no embedding, no coordinate change---the exchange
constant is unnecessary.  The Level~1 $S^3$ graph has integer
eigenvalues; $\kappa$ is needed only to convert them to Hartrees.
The Level~3 FCI on $S^3 \times S^3$ has algebraic matrix elements;
$\mu(R)$ appears only in the hyperspherical reparameterization.
The Level~2 $\sigma$-state solver has algebraic Laguerre moments;
$e^a E_1(a)$ appears only for $m \neq 0$ states where the
coordinate curvature introduces a pole.

This suggests a design principle for the GeoVac program:
\textit{minimize the number of coordinate projections} to minimize
the number of exchange constants.  The ideal solver stays on the
graph; every departure from the graph incurs a transcendental
toll.

\subsection{Exchange constants are computationally cheap}

A striking feature of the exchange constants is that their
conceptual complexity does not correlate with computational cost:
\begin{itemize}
  \item $\kappa = -1/16$: one multiplication.
  \item $e^a E_1(a)$: one function evaluation, followed by an
    $O(N^2)$ algebraic recurrence.
  \item $\mu(R)$: a $10$--$30$ dimensional matrix diagonalization
    per $R$-point (milliseconds).
\end{itemize}
The hierarchy of transcendence---constant, seed, function---tracks
the severity of curvature mismatch, not computational difficulty.
This reinforces the interpretation that exchange constants encode
\textit{geometric information} (the shape of the mismatch between
discrete and continuous descriptions) rather than
\textit{computational obstacles}.  The Level~2 result is
particularly instructive: identifying the Stieltjes seed and
absorbing it into the basis actually \textit{reduces} the
computational cost (by eliminating quadrature) while simultaneously
providing better convergence (Sec.~\ref{sec:stieltjes}).

\subsection{The hierarchy of transcendence}

The exchange constants increase in complexity as the curvature
mismatch increases:
\begin{itemize}
  \item Constant $\to$ constant (Level 1): $\kappa = -1/16$.
    One projection, fixed curvature.
  \item Constant $\to$ single transcendental (Level 2, $m \neq 0$):
    $e^a E_1(a)$.  One embedding with a coordinate singularity.
  \item Constant $\to$ algebraic function of one parameter
    (Level 3, hyperspherical): $\mu(R)$.  Defined implicitly by
    $P(R,\mu) = 0$, a polynomial with coefficients in
    $\mathbb{Q}(\pi, \sqrt{2})$.  The transcendental content is
    inherited from the angular coupling matrix, not generated by
    the parameterization.  Branch points of the algebraic curve
    (eigenvalue crossings in the complex $R$-plane) set the
    convergence radius of perturbation series expansions.
  \item Constant $\to$ piecewise-algebraic function of two parameters
    (Level 4, mol-frame hyperspherical): $\mu(\rho, R)$.
    The split-region Legendre expansion introduces a region boundary
    at $\rho = 1/2$ that breaks the linear pencil structure, so
    $\mu(\rho)$ does not satisfy a global polynomial equation.
    Within each region, the parameter dependence is algebraic.
    The hyperspherical parameterization remains computationally
    essential: the product-space CI saturates at 92.5\%
    $D_e$~\cite{loutey_track_m}.
\end{itemize}

This hierarchy mirrors the mathematical classification of constants
and functions: the Weyl constants are periods
($\pi = \int dx/(1+x^2)$); the Stieltjes seeds are exponential
periods ($E_1(a) = \int_1^\infty e^{-at}/t\,dt$); and the adiabatic
eigenvalue functions are algebraic functions of the continuous
parameters, built from transcendental data.  The ``non-period
transcendental'' classification previously applied to $\mu(R)$ is
therefore superseded: $\mu(R)$ is not transcendental as a function
of $R$, though its coefficients are transcendental numbers.
At Level~4, the piecewise-smooth parameter dependence places $\mu(\rho)$
outside even the algebraic category, though it remains piecewise-algebraic
within each Legendre region.


\subsection{Compactness at the construction level}

The compactness thesis applies not only to the exchange constants but
to the packing construction itself.  Paper~0's Axiom~2 (geometric
indifference) forces states onto isotropic shells---copies of $S^1$,
the simplest compact manifold.  The circumference of the circle
introduces $\pi$ into the fundamental area $\sigma_0 = \pi d_0^2/2$
and thus into every state count $N_k = A_k/\sigma_0 = 2(2k-1)$.
The graph is $\pi$-free in its algebraic operations
(Observation~\ref{obs:pi-free}), but the node count per shell---$2l+1$
for angular momentum $l$---reflects the compact $S^1$ topology of the
angular shell.  The $\pi$ in $\sigma_0$ is not a coordinate artifact; it
is the conversion factor between the discrete count of states fitting on
a compact circle and the continuous Euclidean area that circle encloses.
This is the exchange constant pattern at its most elementary: compact
geometry determines discrete quantum numbers, and $\pi$ is the toll for
expressing them in terms of continuous measures.

\subsection{Structural parallels}

The pattern of compact geometry determining coupling constants via
projection has a structural parallel in string compactification, where
the geometry of a compact internal manifold determines the low-energy
coupling constants of the effective theory through dimensional
reduction~\cite{polchinski1998}.  In both cases, the compact manifold's
topology constrains the discrete spectrum, and the coupling constants
encode the projection from the compact description to the non-compact
one in which measurements are performed.  This is a mathematical
observation about the structure of the projection, not a physical claim
about the relationship between the two frameworks.  We note also that
$\mathrm{SO}(4,2) = \mathrm{Iso}(\mathrm{AdS}_5)$, the isometry group
of five-dimensional anti-de~Sitter space, is the conformal group of the
boundary $\mathbb{R}^{3,1}$; the Fock projection
$S^3 \to \mathbb{R}^3$ is a Euclidean section of this conformal
boundary correspondence.

\subsection{The algebraic curve refinement}
\label{sec:algebraic_curve}

The characteristic polynomial analysis~\cite{loutey_track_p1}
refines the exchange constant taxonomy by distinguishing between
transcendental \emph{constants} and algebraic \emph{functions} of
continuous parameters:
\begin{itemize}
  \item \textbf{Truly irreducible transcendentals:}
    The embedding constants $e^a E_1(a)$ (Level~2, $m \neq 0$) and the
    angular coupling integrals $c/\pi$ and $c\sqrt{2}/\pi$
    (Level~3 and~4) are irreducible---they cannot be eliminated by
    any algebraic manipulation.  These are the genuine ``exchange
    constants'' in the Weyl--Selberg sense.
  \item \textbf{Algebraic functions of transcendental data:}
    The adiabatic eigenvalues $\mu(R)$ and $\mu(\rho,R)$ are
    algebraic functions whose coefficients inherit transcendentals
    from the coupling matrix.  The $R$-dependence (or $(\rho,R)$-dependence)
    is algebraic at every truncation: $\mu$ satisfies a polynomial
    equation $P(R,\mu) = 0$ of degree $d$ equal to the matrix dimension.
    At $l_{\max} = 0$, $P$ is linear and $\mu(R)$ is an explicit
    rational function of $R$ (with transcendental coefficients).
    At $l_{\max} = L$, $P$ has degree $L+1$ in both $R$ and $\mu$.
  \item \textbf{Calibration constant:} $\kappa = -1/16$ is rational
    (Sec.~\ref{sec:kappa_derivation}).
\end{itemize}

This refinement has a practical implication: the ``transcendental toll''
of the adiabatic parameterization is \emph{inherited}, not \emph{accrued}.
Moving from the product space $S^3 \times S^3$ to hyperspherical
coordinates does not introduce new transcendental content; it
rearranges existing transcendentals (from the angular coupling)
into an algebraic function of $R$.  The toll is paid once, in the
coupling matrix, and the $R$-dependence is algebraically determined.

This does not yet eliminate the computational cost of pointwise
diagonalization: at $l_{\max} \geq 1$, the characteristic polynomial
has degree $\geq 2$ in $\mu$ and must be solved numerically at each
$R$-point.  But it does mean that the solver operates on a known
polynomial, not an opaque transcendental function.  Whether this
algebraic structure can be exploited computationally---for instance,
via resultant methods, homotopy continuation, or algebraic
PES representations---is an open question.

This algebraic structure is specific to Level~3, where the angular
Hamiltonian $H(R) = H_0 + R \cdot V_C$ is a linear matrix pencil.
At Level~4, the split-region Legendre expansion of the nuclear charge
function introduces piecewise-smooth $\rho$-dependence that breaks
the pencil structure (Track~S).  The Level~4 angular eigenvalues
are piecewise-algebraic but do not satisfy a global characteristic
polynomial.

\subsection{Derivation of $\kappa = -1/16$}
\label{sec:kappa_derivation}

The original identification of $\kappa = -1/16$ (Paper~0) was
empirical: the unique constant that maps GeoVac graph eigenvalues
onto the Rydberg energy levels.  We now derive this value and
classify its origin.

\textit{Graph Laplacian spectrum.}  The GeoVac lattice
(Paper~0) discretizes $S^3$ with states $|n,l,m\rangle$
connected by angular transitions $m \leftrightarrow m \pm 1$
and radial transitions $n \leftrightarrow n \pm 1$.  Interior
nodes have degree $d_{\max} = 4$ (two angular, two radial
neighbors).  The graph Laplacian $\mathcal{L} = D - A$ is
positive semidefinite; its maximum eigenvalue satisfies
$\lambda_{\max} \leq 2 d_{\max} = 8$, with the bound saturated
in the continuum limit ($n_{\max} \to \infty$).  Numerical
verification confirms monotonic convergence:
$\lambda_{\max}(n_{\max}{=}5) = 6.62$,
$\lambda_{\max}(10) = 7.61$,
$\lambda_{\max}(20) = 7.90 \to 8$.

\textit{Energy calibration.}  The graph Hamiltonian
$H = \kappa \mathcal{L}$ (no node weights) has its ground
state at energy $\kappa \lambda_{\max}$, since $\kappa < 0$
maps the largest graph eigenvalue to the lowest physical energy.
Matching this to the hydrogen ground state:
\begin{equation}
  \kappa \lambda_{\max} = E_H = -\tfrac{1}{2}
  \implies
  \kappa = \frac{-1/2}{8} = -\frac{1}{16}.
  \label{eq:kappa_derived}
\end{equation}

\textit{Factor identification.}  The value $-1/16 = -(1/2)(1/8)$
decomposes into:
\begin{enumerate}
  \item $1/2$: the non-relativistic kinetic energy normalization
    $E = p^2/(2m)$, or equivalently the energy-shell constraint
    $E = -p_0^2/2$;
  \item $1/8 = 1/(2 d_{\max})$: the inverse of the maximum
    graph Laplacian eigenvalue, determined by the vertex degree
    of the lattice discretization.
\end{enumerate}

\textit{Heat kernel test.}  The Seeley--DeWitt expansion of the
heat kernel on unit $S^3$ gives coefficients
$a_0 = \mathrm{vol}(S^3) = 2\pi^2$,
$a_1 = (R/6)\,\mathrm{vol}(S^3) = 2\pi^2$,
$a_2 = \frac{1}{360}(5R^2 - 2|\mathrm{Ric}|^2 +
2|\mathrm{Riem}|^2)\,\mathrm{vol}(S^3) = \pi^2$,
using scalar curvature $R = 6$, $|\mathrm{Ric}|^2 = 12$,
$|\mathrm{Riem}|^2 = 12$ for constant sectional curvature
$K = 1$.  No ratio of these coefficients produces $1/16$.  The
Weyl coefficient is $N(\Lambda) \sim \Lambda^{3/2}/3$, also
unrelated.  The spectral zeta function
$\zeta_{S^3}(s) = \sum_{l \geq 1} (l+1)^2/[l(l+2)]^s$
similarly produces no value matching $1/16$.

\textit{Conformal coupling coincidence.}  The conformal
coupling constant for a 3-manifold is $c_3 = (d-2)/[4(d-1)]
= 1/8$ for $d = 3$, numerically equal to $1/\lambda_{\max}$.
This coincidence is not causal: $c_3$ arises from the Yamabe
operator $Y = -\Delta + (R/8)$ on $S^3$, while
$1/\lambda_{\max}$ arises from the maximum vertex degree of
the graph.  A different discretization of $S^3$ (e.g., adding
$l \leftrightarrow l \pm 1$ transitions to reach $d_{\max} = 6$)
would give $\lambda_{\max} \to 12$ and $\kappa \to -1/24$,
breaking the coincidence while preserving $c_3 = 1/8$.

\textit{Classification.}  $\kappa$ is correctly classified as
``conformal'' in Table~\ref{tab:taxonomy}: it is needed because
the Fock stereographic projection introduces a conformal factor
relating the $S^3$ Laplacian to flat-space operators.  However,
its \textit{numerical value} is a graph calibration constant
determined by the lattice connectivity, not an intrinsic
spectral invariant of $S^3$.  The Z-universality noted in
Paper~0 ($\kappa \to \kappa Z^2$ for any nuclear charge)
reflects the Z-independence of the graph structure, with only
the energy-shell scaling $E \propto Z^2$.


%% ========================================================================
%% VIII. CONCLUSION
%% ========================================================================

\section{Conclusion}
\label{sec:conclusion}

Compactness is the organizing principle behind the exchange constants
cataloged in this paper.  Compact Lie groups have discrete unitary
representations (Peter--Weyl); the GeoVac graph exploits this
discreteness at every level.  Transcendental conversion factors
arise when---and only when---those discrete representations are
projected onto non-compact coordinate systems for physical
interpretation.  The taxonomy classifies these factors by the type
of compactness mismatch: from the elementary $\pi$ of the $S^1$
packing shell, through the calibration $\kappa$ of the $S^3 \to
\mathbb{R}^3$ stereographic projection, to the composition constants
that measure the cost of factorizing a single compact space into
multiple coordinate systems.  The irreducible transcendentals---$\pi$,
$e^a E_1(a)$---are instances of the Weyl--Selberg exchange
constants, encoding the curvature mismatch between the graph and
the coordinate system.  The Level~3 adiabatic eigenvalue $\mu(R)$,
previously classified as transcendental, is in fact an algebraic
function of $R$ with transcendental coefficients inherited from the
angular coupling matrix.  The Level~4 eigenvalue $\mu(\rho,R)$ is
piecewise-algebraic, with the region structure inherited from the
split-region Legendre expansion.  The calibration constant
$\kappa = -1/16$ is rational.  All exchange constants vanish when
the computation stays on the graph.

At Level~4, the situation is qualitatively different from
Levels~1--3: the product-space computation (prolate spheroidal
CI) is fully algebraic but saturates at 92.5\% of the exact
dissociation energy (Track~M~\cite{loutey_track_m}).  The
remaining 7.5\%---the electron-electron cusp---requires the
flow function $\mu(\rho,R)$, which is not merely a projection
toll but an irreducible encoding of three-body coalescence
physics that the algebraic product space cannot represent.

The observation that the conjectural fine structure constant
formula $K = \pi(B + F - \Delta)$ has the same structural
character---a Weyl-type conversion factor mixing discrete
representation data with continuous spectral invariants---suggests
that the combination rule may be derivable from the spectral
geometry of the Hopf bundle, rather than being an irreducible
empirical observation.

Two implications for the GeoVac program follow:
\begin{enumerate}
  \item \textbf{The design principle}: Staying on the graph
    is not merely a computational preference; it is the
    algebraically natural thing to do.  Every coordinate
    projection incurs a transcendental toll, though that toll
    may be a constant (inherited from angular coupling) rather
    than a function.  The adiabatic parameterization introduces
    no new transcendental content; it rearranges upstream constants
    into an algebraic function of $R$.  The program's ongoing effort
    to algebraicize the solvers (Paper~12's Neumann expansion,
    the Laguerre moment recurrences, the Gaunt integral replacements)
    is, in retrospect, an effort to eliminate exchange constants by
    eliminating projections---and the algebraic curve result shows
    that this effort has been more successful than previously
    recognized.  The cross-center nuclear attraction integrals of
    Paper~19 provide the most recent example: the multipole expansion
    of $V_{ne}^{(b,a)}$ terminates exactly at $L_{\max} = 2 l_{\max}$
    by Gaunt selection rules, and the radial split integrals evaluate
    analytically via incomplete gamma functions for hydrogenic
    wavefunctions (machine precision, zero quadrature).  In the
    exchange constant taxonomy, the grid-based $V_{ne}$ was a
    calibration exchange constant---transcendental content from
    projection between the single-center basis and the two-center
    geometry---that proved algebraic once the polynomial-times-exponential
    structure of the integrand was identified.
  \item \textbf{The physical interpretation of $\kappa$}:
    Paper~0 identified $\kappa = -1/16$ as a universal packing
    constant but did not derive it from first principles.
    Explicit computation (Sec.~\ref{sec:kappa_derivation})
    shows that $\kappa$ is \textit{not} computable from the
    heat kernel expansion of the $S^3$ Laplacian: the
    Seeley--DeWitt coefficients $a_0 = 2\pi^2$, $a_1 = 2\pi^2$,
    $a_2 = \pi^2$ do not produce $1/16$ in any natural ratio.
    Instead, $\kappa = -E_H / \lambda_{\max}(\mathcal{L})$,
    where $E_H = -1/2$~Ha is the hydrogen ground state energy
    and $\lambda_{\max} \to 8$ is the maximum eigenvalue of
    the GeoVac graph Laplacian in the continuum limit
    ($8 = 2 \times d_{\max}$, where $d_{\max} = 4$ is the
    maximum vertex degree).  The Table~\ref{tab:taxonomy}
    classification as ``conformal'' remains appropriate---$\kappa$
    is needed because of the conformal projection---but its
    numerical value depends on the graph construction, not on
    abstract spectral invariants of $S^3$.
\end{enumerate}

\paragraph{Note added v2.0.24.}
The exchange constant taxonomy has been validated by the full
$N$-electron solver (Tracks~AJ--AS, v2.0.19--23).  The PK
pseudopotential, classified here as a \emph{composition exchange
constant}, has been confirmed as an approximation to exact $S_4$
antisymmetry: the full 4-electron solver produces molecular
equilibrium without PK at $l_{\max} \geq 2$ (Paper~17, Sec.~VIII).
The 2D variational solver (Track~AR) confirmed that adiabatic $D_e$
overcounting was an artifact, and that the $l_{\max} = 2$ angular
basis is genuinely insufficient for binding---the composed geometry's
$144{\times}$ angular compression via core/valence pre-optimization
is essential.  The quantum encoding comparison (Track~AS) confirmed
that the composed architecture produces categorically sparser qubit
Hamiltonians than the full $N$-electron encoding, with the structural
sparsity arising from Gaunt selection rules (the same algebraic
structure that produces the exchange constants classified in this
paper).  The classical solver investigation across all levels is now
complete (v2.0.24).

\paragraph{Note added v2.6.0.}
The He exchange constant census is now computationally explicit
(Track~DI Sprints~1--3D).  A 2D variational solver on $S^5$ breaks
the adiabatic 0.20\% floor: 0.022\% raw at $l_{\max} = 7$, 0.004\%
with Schwartz cusp correction at $l_{\max} = 4$ (Paper~13, Note
added v2.6.0).  A graph-native FCI using the graph Laplacian
$\kappa\mathcal{L}$ as the one-body operator and exact rational
Slater integrals as the two-body operator achieves 0.19\% at
$n_{\max} = 7$ (1,218~configurations) with zero free parameters.
The graph's off-diagonal $T_\pm$ inter-shell couplings are the
dominant ingredient: replacing $\kappa\mathcal{L}$ with the diagonal
hydrogenic operator degrades to 1.41\% (variational~$k$) or 1.85\%
(fixed $k = Z$).  FCI basis invariance verified (Sprint~3D):
transforming $V_{ee}$ to the graph eigenbasis gives identical
energies ($< 3 \times 10^{-15}$~Ha), confirming the $\sim$0.14\%
convergence floor is embedding exchange constant content (cusp),
not basis mismatch.  The three-layer structure is refined:

\begin{center}
\begin{tabular}{llll}
\hline
Layer & Content & Exchange type & Contribution \\
\hline
Rational & Graph h$_1$ + Slater $V_{ee}$ & None (graph-intrinsic) & 98.6\% \\
Topological & Graph off-diagonal ($T_\pm$) & Calibration ($\kappa$) & Dominant \\
Transcendental & e-e cusp + truncation & Embedding & $\sim$0.14\% \\
\hline
\end{tabular}
\end{center}

\noindent
The calibration exchange constant $\kappa = -1/16$ has a dual role:
it sets the overall energy scale (Level~1) and determines the
strength of inter-shell coupling in the graph Laplacian (Level~3).
The Fock self-consistency $k^2 = -2E$, exact at Level~1, fails at
Level~3 (5--13\% error), marking the transition from single
calibration constants to irreducible embedding content.

\begin{acknowledgments}
Computational implementation and documentation drafting were
performed with the assistance of large language models (Anthropic
Claude) under the author's scientific direction.  All results
were validated against analytical benchmarks.
\end{acknowledgments}


%% ========================================================================
%% BIBLIOGRAPHY
%% ========================================================================

\begin{thebibliography}{30}

\bibitem{loutey_paper0}
J.~Loutey,
``The Geometric Atom: Quantum State Space as a Packing Problem,''
GeoVac Paper~0 (2026).

\bibitem{loutey_paper1}
J.~Loutey,
``The Geometric Atom: Quantum Mechanics as a Packing Problem,''
GeoVac Paper~1 (2026).

\bibitem{loutey_paper7}
J.~Loutey,
``The Dimensionless Vacuum: Recovering the Schr\"odinger Equation
from Scale-Invariant Graph Topology,''
GeoVac Paper~7 (2026).

\bibitem{loutey_paper11}
J.~Loutey,
``The Molecular Fock Projection: Diatomic Molecules as Prolate
Spheroidal Lattices,''
GeoVac Paper~11 (2026).

\bibitem{loutey_paper13}
J.~Loutey,
``The Hyperspherical Lattice: Two-Electron Atoms as Coupled
Channel Graphs,''
GeoVac Paper~13 (2026).

\bibitem{loutey_fci}
J.~Loutey,
``Full configuration interaction on a sparse graph Hamiltonian:
Multi-electron atoms via relational lattice indexing,''
GeoVac FCI Paper (2026).

\bibitem{loutey_paper2}
J.~Loutey,
``The Fine Structure Constant from Spectral Geometry of the Hopf
Fibration,''
GeoVac Paper~2 (2026).

\bibitem{loutey_track_j}
J.~Loutey,
``Track J: Algebraic Level 2 Radial Solver for $m \neq 0$ States,''
GeoVac supplementary material (2026).

\bibitem{loutey_track_g}
J.~Loutey,
``Track G: Perturbation Series for $\mu(R)$,''
GeoVac supplementary material (2026).

\bibitem{loutey_track_m}
J.~Loutey,
``Track M: Prolate CI $l_{\max}$ Convergence Study,''
GeoVac supplementary material (2026).

\bibitem{loutey_track_p1}
J.~Loutey,
``Track P1: Algebraic Curve for $\mu(R)$---Characteristic Polynomial
of the Level~3 Angular Hamiltonian,''
GeoVac supplementary material (2026).

\bibitem{loutey_track_s}
J.~Loutey,
``Track S: Level~4 Algebraic Curve---Structural Negative Result,''
GeoVac supplementary material (2026).

\bibitem{drake2006}
G.~W.~F.~Drake,
``High precision calculations for helium,''
in \textit{Springer Handbook of Atomic, Molecular, and Optical Physics},
edited by G.~W.~F.~Drake (Springer, New York, 2006), pp.~199--219.

\bibitem{dlmf_laguerre}
NIST Digital Library of Mathematical Functions,
\S18.9.13, \url{https://dlmf.nist.gov/18.9.13}.

\bibitem{weyl1911}
H.~Weyl,
``\"Uber die asymptotische Verteilung der Eigenwerte,''
Nachr.\ K\"onigl.\ Ges.\ Wiss.\ G\"ottingen, Math.-phys.\ Kl.\
(1911), pp.\ 110--117.

\bibitem{selberg1956}
A.~Selberg,
``Harmonic analysis and discontinuous groups in weakly symmetric
Riemannian spaces with applications to Dirichlet series,''
J.\ Indian Math.\ Soc.\ \textbf{20}, 47--87 (1956).

\bibitem{berger2003}
M.~Berger,
\textit{A Panoramic View of Riemannian Geometry}
(Springer, Berlin, 2003).

\bibitem{lichnerowicz1963}
A.~Lichnerowicz,
\textit{Spineurs harmoniques},
C.~R.\ Acad.\ Sci.\ Paris \textbf{257}, 7--9 (1963).

\bibitem{camporesi_higuchi1996}
R.~Camporesi and A.~Higuchi,
On the eigenfunctions of the Dirac operator on spheres and
real hyperbolic spaces,
\textit{J.\ Geom.\ Phys.}\ \textbf{20}, 1--18 (1996).

\bibitem{peter_weyl1927}
F.~Peter and H.~Weyl,
``Die Vollst\"andigkeit der primitiven Darstellungen einer
geschlossenen kontinuierlichen Gruppe,''
\textit{Math.\ Ann.}\ \textbf{97}, 737--755 (1927).

\bibitem{polchinski1998}
J.~Polchinski,
\textit{String Theory}, Vol.~2: Superstring Theory and Beyond
(Cambridge University Press, Cambridge, 1998).

\bibitem{goncharov1999}
A.~B.~Goncharov,
``Volumes of hyperbolic manifolds and mixed Tate motives,''
\textit{J. Amer. Math. Soc.} \textbf{12}, 569--618 (1999).

\bibitem{zagier1994}
D.~Zagier,
``Values of zeta functions and their applications,''
in \textit{First European Congress of Mathematics}, Vol.~II
(Birkh\"auser, Basel, 1994), pp.~497--512.

\end{thebibliography}

\end{document}
